\section{Methods}\label{sec:methods}

This section develops the mathematical core of the framework, our approach to the frameworks of \S\ref{sec:manin}: Manin's quantum tori $\F_1$ program with Borger's $\Lambda$-algebraic ring. Here we develop only what is pertinent to the string; not the whole of the $\F_1$ program, which is developed in a companion paper~\cite{F1}. What we develop here is presented as a single thread, from $i$ to the two vertices $\mu=\tfrac12$ and $\sigma=\tfrac32$,
each step formally verified in Lean~4 (theorem names in brackets); each
subsection uses only objects already introduced.

% ---------------------------------------------------------------------
\noindent\emph{Verification tags:} $\P{...}$ formal proof (Lean), $\N{...}$ numerical check,
$\Hyp{...}$ hypothesis from the literature, $\Conj$ proposed. Compound tags
($\P{...}\;\N{...}$) mark a claim resting on more than one, and each input is named where it enters.

\subsection{The generator, its conjugate, and the binary bridge}\label{ssec:generator}

Everything below is read off one quadratic. We fix here what $\golden$ is, what its conjugate
is, and the one identity of \S\ref{sec:intro} we shall use throughout; nothing in this
subsection depends on anything later.

\begin{equation}\label{eq:base}
\golden^2=\golden+1,\qquad
\Rpcf=\Z[\golden,\golden^{-1},\tfrac12],\qquad
\E^3=\{(x,y,\golden y)\}.
\end{equation}

\begin{definition}[Conjugate, trace, norm, discriminant]\label{def:K-arith}
Write $\bar\golden=(1-\sqrt5)/2$ for the Galois conjugate of $\golden$, the second root of
$x^{2}-x-1$. Then
\begin{equation}\label{eq:trace-norm}
  \golden+\bar\golden=1,
  \qquad
  \golden\,\bar\golden=-1,
  \qquad
  (\golden-\bar\golden)^{2}=\Delta_K=5 .
\end{equation}
The norm $N(\golden)=-1$ exhibits $\golden$ as a unit; the discriminant $\Delta_K=5$ makes $5$
the only ramified rational prime, which is what the trichotomy of \S\ref{ssec:spectrum} rests on.
\P{phi_trace, phi_norm, galois_pair, phi_discriminant}
\end{definition}

\begin{lemma}[The Galois conjugate is the involution]\label{lem:galois-inv}
$\bar\golden=1-\golden$; equivalently, $x\mapsto(\golden+\bar\golden)-x$ and $x\mapsto1-x$ are
the same map. The Galois conjugation of $K$ and the involution $x\mapsto1-x$ are therefore not
two analogous symmetries but one map, read on the orbit $\{\golden,\bar\golden\}$ and on $\R$.
\P{galois_conj_is_one_sub, galois_involution_is_one_sub}
\end{lemma}

Finally we record, without reproving it, the identity established at \eqref{eq:bridge}:
$\golden^{\lng}=2$ with $\lng=\log_{\golden}2$. It is the only result of \S\ref{sec:intro} that
the present section consumes. \P{mersenne\_bridge}

\begin{proposition}[The bridge anchors from outside the lattice]\label{prop:anchor-exterior}
The exponent $\lng$ is not reachable by descent on the lattice of integer exponents. For no integer
$n$ does $\golden^{n}=2$:
\begin{equation}\label{eq:anchor-exterior}
  \golden<2<\golden^{2}=\golden+1,
\end{equation}
and $n\mapsto\golden^{n}$ is strictly monotone since $\golden>1$, so by \eqref{eq:anchor-exterior}
the $2$ falls in the gap between two consecutive powers and is missed by all of them. The identity
$\golden^{\lng}=2$ therefore holds in $\R$ and at no point of $\Z$. This is what makes
\eqref{eq:bridge} an \emph{anchor}: it fixes the scale from outside rather than being derived
inside the tower, which is why the present section consumes it and does not prove it.
\P{phi_lt_two, two_lt_phi_sq, two_not_in_golden_lattice, anchor_is_exterior}
\end{proposition}

% ---------------------------------------------------------------------
\subsection{Three geometric origins of $\mu=\tfrac12$}\label{ssec:origins}

The constant $\mu=\tfrac12$ arises from three independent origins on the generator of
\S\ref{ssec:generator}: a $\pi$-branch (Gamma/Gaussian), a $\golden$-branch (the involution
$x=1-x$), and an arithmetic branch (the Galois pair of $K$). These are the seeds.

\begin{theorem}[Pentagonal identity]\label{thm:pentagon-id}
$\golden=2\cos(\pi/5)$. \P{phi_eq_two_cos_pi_fifth}
\end{theorem}
\medskip
\begin{proof}
The Chebyshev polynomial $T_5$ satisfies $T_5(\cos\theta)=\cos(5\theta)$ with
$T_5(x)=16x^5-20x^3+5x$. At $\theta=\pi/5$ this gives
$T_5(\cos(\pi/5))=\cos\pi=-1$, so $\cos(\pi/5)$ is a root of
$16x^5-20x^3+5x+1=0$, which factors as $(x+1)(4x^2-2x-1)^2$. Since
$0<\pi/5<\pi/2$ forces $\cos(\pi/5)>0$, we have $\cos(\pi/5)\neq-1$; hence
$\cos(\pi/5)$ is a root of $4x^2-2x-1=0$, whose positive root is
$(1+\sqrt5)/4=\golden/2$. Therefore $\golden=2\cos(\pi/5)$; this classical identity is also
derived in the correspondence paper~\cite{Corr}.
\end{proof}

\begin{proposition}[Pentagonal bridge $\pi\leftrightarrow\golden$]\label{prop:pi-bridge}
Since $\golden=2\cos(\pi/5)$ (Theorem~\ref{thm:pentagon-id}),
\begin{equation}\label{eq:pi-arccos}
\pi = 5\arccos(\golden/2).
\end{equation}
\P{}
\end{proposition}

\begin{proposition}[Gamma as an infinite integral (Euler)]\label{prop:gamma-int}
For $s>0$,
\begin{equation}\label{eq:gamma-int}
\Gamma(s)=\int_0^\infty t^{\,s-1}e^{-t}\,dt .
\end{equation}
\P{}
\end{proposition}

\begin{lemma}[Gaussian value]\label{lem:gamma-half}
$\Gamma(\tfrac12)=\sqrt\pi$ (classical; from the Gaussian integral $\int_{-\infty}^{\infty}e^{-x^2}dx=\sqrt\pi$).
\P{}
\end{lemma}
\medskip
\begin{theorem}[Half-factorial (classical)]\label{thm:half-factorial}
\begin{equation}\label{eq:half-factorial}
(\tfrac12)! \;=\; \Gamma(\tfrac32) \;=\; \mu\sqrt\pi,
\qquad\text{equivalently}\qquad \mu=\frac{(\tfrac12)!}{\sqrt\pi}.
\end{equation}
\P{}
\end{theorem}
The Gaussian of Lemma~\ref{lem:gamma-half} enters with its width still free; three separate
conditions close it.

\begin{proposition}[Uniqueness of the self-dual Gaussian]\label{prop:selfdual-gaussian}
Let $g_a(x)=e^{-a x^{2}}$ be the Gaussian family of which
Lemma~\ref{lem:gamma-half} uses the case $a=1$, with the width $a>0$ now left free, and let $\widehat g(\xi)=\int_{\R} g(x)e^{-2\pi i x\xi}\,dx$, so that
$\widehat{g_a}(\xi)=\sqrt{\pi/a}\;e^{-\pi^{2}\xi^{2}/a}$. Three independent conditions
each fix $a=\pi$: normalisation, the amplitude of the Fourier
transform, and its width. The self-dual Gaussian is therefore
\begin{equation}\label{eq:selfdual-gaussian}
  g(x)=e^{-\pi x^{2}},\qquad \widehat g=g,\qquad \int_{\R}g=1 .
\end{equation}
\P{gauss_normalised_iff, fourier_amplitude_iff, fourier_width_iff, selfdual_gaussian_unique}
\end{proposition}
\begin{proof}
$\int_{\R}g_a=\sqrt{\pi/a}$, so normalisation gives $\sqrt{\pi/a}=1$, hence $a=\pi$.
Self-duality $\widehat{g_a}=g_a$ gives independently $\sqrt{\pi/a}=1$ (amplitude) and
$\pi^{2}/a=a$ (width); since $a,\pi>0$, the second reads $(a-\pi)(a+\pi)=0$ with
$a+\pi>0$, hence again $a=\pi$. The three conditions are distinct and coincide.
\end{proof}


\begin{proposition}[Mediation half-powers]\label{prop:mediation}
\begin{equation}\label{eq:mediation}
\golden^{\mu\lng}=\sqrt2,\qquad
\golden^{\mu\log_\golden 3}=\sqrt3 .
\end{equation}
\P{}
\end{proposition}

\begin{proposition}[$\golden$-branch: the involution and its unique fixed point]\label{prop:phi-branch}
$\mu=\tfrac12$ is the \emph{unique} fixed point of $x\mapsto1-x$: for real $x$, $x=1-x$ if and
only if $x=\tfrac12$. \P{faceInv_apex}
\end{proposition}

\begin{proposition}[Arithmetic branch: the midpoint of the Galois pair]\label{prop:galois-seed}
By Lemma~\ref{lem:galois-inv} the Galois conjugation of $K$ and the involution $x\mapsto1-x$ are
the same map; the midpoint of the pair $\{\golden,\bar\golden\}$ is therefore its fixed point,
\begin{equation}\label{eq:galois-seed}
  \frac{\golden+\bar\golden}{2}=\frac12=\mu ,
\end{equation}
by the trace of \eqref{eq:trace-norm}. This is the third seed, and it is the one that is
\emph{arithmetic}: the other two are analytic and elementary, whereas this one is an invariant
of the field $K$ itself.
\P{faceGalois_apex}
\end{proposition}

% ---------------------------------------------------------------------
\subsection{Arity uniqueness and the blocking of the diagonal}\label{ssec:arity}

The arity $3$ and the modulus $\tfrac12$ are forced, and
the Lawvere--Yanofsky diagonal is contracted (not prohibited).

\begin{definition}[Distributed self-reference]\label{def:dsr}
A recurrence of depth $k$ is \emph{distributed} when $k\ge 2$ (no component
depends directly on itself). \P{DistributedSelfReference}
\P{}
\end{definition}

\begin{theorem}[Fibonacci minimality]\label{thm:fib-min}
Depth $k=2$ is minimal, and the unique positive root of $r^2=r+1$ is
$\golden$. Among natural-coefficient depth-$2$ recurrences
$F_{n+2}=c_1F_{n+1}+c_2F_n$, the pair $c_1=c_2=1$ is the unique one whose
characteristic root---hence growth ratio---is $\golden$.
\P{} \N{}
\end{theorem}

\begin{proposition}[Spectral constraints force $(\sigma,\mu)=(\tfrac32,\tfrac12)$]\label{prop:spectral}
$\sigma+\mu=2$, $\sigma\mu=\tfrac34$, with unique solution
$0<\mu<1$ giving $\mu=\tfrac12$, $\sigma=\tfrac32$.
\P{}
\end{proposition}

\begin{remark}[No-diagonal]\label{rmk:no-diagonal}
$\|\hat\Omega\|=\tfrac12<1$ contracts the diagonal below the contradiction
threshold.
\end{remark}

% ---------------------------------------------------------------------
\subsection{Construction of PCF and the base ring on the torus}\label{ssec:construction}

On the generator and the conjugate of \S\ref{ssec:generator}, the ring of integers and the
regulator, then the base ring, the tripartite operator, and the projection that fixes
$\varepsilon_0$ out of the regulator.


\begin{proposition}[The ring of integers and the base ring]\label{prop:rings}
\begin{equation}\label{eq:OK-vs-Rpcf}
  \OK=\Z[\golden],
  \qquad
  \Rpcf=\Z\bigl[\golden,\golden^{-1},\tfrac12\bigr] .
\end{equation}
The two differ by integrality: $\golden$ is a root of the \emph{monic} $x^{2}-x-1$, so it is an
algebraic integer and lies in $\OK$, whereas $\tfrac12$ is a root of $2x-1$, which is not monic.
$\Rpcf$ therefore adjoins to $\OK$ an element that is not an algebraic integer, namely
$\tfrac12$; the modulus that the tripartite operator will carry is that same $\tfrac12$, a unit
of $\Rpcf$ and not of $\OK$. \P{phi_is_algebraic_integer, half_not_algebraic_integer}

\end{proposition}

\begin{definition}[Regulator and period]\label{def:regulator}
\begin{equation}\label{eq:regulator}
  R_K=\log\golden,
  \qquad
  T=2\pi\log\golden=2\pi R_K .
\end{equation}
$R_K$ is the regulator of $K$ in the arithmetic sense---not the ultraviolet regulators of
\S\ref{ssec:p-tower}---and $T$ is the period of the torus constructed below in this subsection.
Both are named here
because $\varepsilon_0=\ln\golden/(6\sqrt3)$ is derived below out of $\ln\golden$, and the reader
should have the invariant named before it is used. \P{regulator_K_pos, period_eq_two_pi_regulator}

\end{definition}

\begin{equation}\label{eq:omega-norms}
\Omega=P\cdot C\cdot F .
\end{equation}
The three component norms are not posited: they are forced by the eigenvalue geometry, and their
product is the modulus (Proposition~\ref{prop:pcf-norms}, \eqref{eq:norms-derived}).

\begin{proposition}[Component norms from the eigenvalue triangle]\label{prop:pcf-norms}
The operator $\hat\Omega$ has eigenvalues $\lambda_k=\tfrac12\omega^{k}$
($k=0,1,2$, $\omega=e^{2\pi i/3}$), inscribing an equilateral triangle in the
circle $|z|=R=\tfrac12$ of $\C$. The component norms are fixed by this geometry:
\begin{equation}\label{eq:norms-derived}
\|C\|=1,\qquad \|F\|=\tfrac{\sqrt3}{2},\qquad \|P\|=\tfrac{1}{\sqrt3},
\end{equation}
whence $\|P\|\,\|C\|\,\|F\|=\tfrac12$, which is also the operator norm $\|\hat\Omega\|$ and the
common modulus $|\lambda_k|$ of each eigenvalue. Three quantities coincide here and are
distinguished in Proposition~\ref{prop:triad-invariants}, where they cease to: the product of
the three \emph{component} norms, the product of the three eigenvalue \emph{moduli}, and the
product of their \emph{real parts}.
\P{normP_eq_tan, normF_eq_cos, Omega_eigenvalues}
\end{proposition}
\begin{proof}
\emph{(C)} $C$ is the real eigenvalue $\lambda_0=\tfrac12$, normalised by the
circumradius $R=|\lambda_0|=\tfrac12$: $\|C\|=|\lambda_0|/R=1$.
\emph{(F)} $F$ is the chord between adjacent eigenvalues,
$\|F\|=|\lambda_1-\lambda_0|=\tfrac12|\omega-1|=\sin(\pi/3)=\tfrac{\sqrt3}{2}$.
\emph{(P)} The fundamental right triangle (centre, midpoint of
$\lambda_0\lambda_1$, vertex $\lambda_0$) has apothem $R\cos(\pi/3)=\tfrac14$ and
half-side $\|F\|/2=\tfrac{\sqrt3}{4}$, so
$\|P\|=\dfrac{R\cos(\pi/3)}{\|F\|/2}=\dfrac{1/4}{\sqrt3/4}=\tfrac{1}{\sqrt3}=\tan(\pi/6)$.
\emph{(Product)} The $\sqrt3$ of $\|F\|$ cancels that of $\|P\|$, leaving $\tfrac12$;
the cancellation is specific to the equilateral ($n{=}3$) triangle.
\end{proof}


\begin{proposition}[Invariants of the eigenvalue triad]\label{prop:triad-invariants}
For the triad $\lambda_k=\tfrac12\omega^{k}$ of Proposition~\ref{prop:pcf-norms}, with
$\golden^{\lng}=2$ \eqref{eq:bridge}:
\begin{equation}\label{eq:triad-re}
  \Re\lambda_k=\golden^{-\lng}\cos\!\Bigl(\frac{2\pi k}{3}\Bigr),\qquad k=0,1,2,
\end{equation}
so the real spectrum is $(\tfrac12,-\tfrac14,-\tfrac14)$ and its sum vanishes. The two products
do \emph{not} agree, and each exponent records a different face of the same
$\golden^{2}=\golden+1$:
\begin{equation}\label{eq:triad-products}
  \prod_{k}\lvert\lambda_k\rvert=2^{-3},
  \qquad
  \prod_{k}\Re\lambda_k=2^{-5}=\golden^{-5\lng}.
\end{equation}
The base $2$ is the binary bridge \eqref{eq:bridge} in both. The exponent $3$ is the arity of the
triad---the same $n{=}3$ whose cancellation gives $\tfrac12$ in
Proposition~\ref{prop:pcf-norms}---and the exponent $5$ is the pentagon of
Theorem~\ref{thm:pentagon-id}. The two identities that generate the framework meet inside one
invariant of the spectrum.
\P{triad_re, triad_trace_zero, triad_two_products, triad_det_is_phi_pow}
\end{proposition}

\begin{lemma}[Isometric embedding of the triad]\label{lem:isometry}
Let $v_k=\tfrac12\omega^{k}$ ($k=0,1,2$) be the eigenvalue triad of
Proposition~\ref{prop:pcf-norms}, and set $\lVert v\rVert^2=\sum_k|v_k|^2$ and
$V_k=v_k/\lVert v\rVert$. Then
\begin{equation}\label{eq:isometry-triad}
  \lVert v\rVert^2=3\cdot\tfrac14=\tfrac34,\qquad \lVert v\rVert=\tfrac{\sqrt3}{2},\qquad
  \sum_{k=0}^{2}\lVert V_k\rVert^2=1,
\end{equation}
so the column map $V\colon\C\to\C^3$, $z\mapsto(V_0,V_1,V_2)\,z$, is norm preserving:
$V^{\dagger}V=\mathbf1$. \P{bulkBoundary\_isometry}
\P{}
\end{lemma}
\begin{proof}
Each $|v_k|=\tfrac12$ (Proposition~\ref{prop:pcf-norms}), so $\lVert v\rVert^2=3\cdot\tfrac14=\tfrac34$
and $\lVert v\rVert=\sqrt3/2$. Then $\lVert V_k\rVert^2=|v_k|^2/\lVert v\rVert^2=(\tfrac14)/(\tfrac34)=\tfrac13$
for each $k$, and the three sum to $1$. For a single column $V$ that sum is $V^{\dagger}V$.
\end{proof}

\begin{remark}[Where the two spectral constraints come from]\label{rmk:spectral-origin}
The system solved in Proposition~\ref{prop:spectral} is not assumed: both of its constraints are
established here. The product is the squared norm of the eigenvalue triad,
$\sigma\mu=\lVert v\rVert^{2}=\tfrac34$~\eqref{eq:isometry-triad}, and the ratio is the arity,
$\sigma/\mu=d=3$ (\S\ref{ssec:arity}). Eliminating $\sigma=3\mu$ gives $3\mu^{2}=\tfrac34$, hence
$\mu=\tfrac12$ and $\sigma=\tfrac32$, and only then $\sigma+\mu=2$. The same $\tfrac34$ reappears
as the colour ratio $1-\muthree^{2}$ and as the GKP entry of the shared signature
(\S\ref{ssec:web}), which is why it is an invariant of the object and not a fitted value.
\P{spectral_uniqueness, bulkBoundary_isometry}
\end{remark}

\begin{definition}[PCF projection]\label{def:projpcf}
For $a,b\in\R$ and $c\neq0$, the \emph{PCF projection} sends the ratio $ab/c$ onto the
$P$-component of the triad and measures it in the angular unit of the eigenvalue triangle:
\begin{equation}\label{eq:projpcf}
  \pi_{\PCF}(a,b,c)\;=\;\frac{ab}{c}\cdot\underbrace{\lVert P\rVert}_{1/\sqrt3}
  \cdot\underbrace{\frac{\pi}{3}}_{\text{angular unit}}
  \;=\;\frac{ab}{c\sqrt3}\cdot\frac{\pi}{3}.
\end{equation}
\P{}
\end{definition}

\begin{remark}[The two constants of \eqref{eq:projpcf} are already fixed]\label{rmk:projpcf-constants}
Neither factor is free. The projection norm is $\lVert P\rVert=1/\sqrt3=\tan(\pi/6)$, forced by the
eigenvalue geometry in Proposition~\ref{prop:pcf-norms}; and $\pi/3$ is the angular unit of the
same equilateral triangle---the half-angle between adjacent eigenvalues $\lambda_k=\tfrac12\omega^k$,
already used as the normalising angle in Proposition~\ref{prop:angle}, where
$\sigma=\zeta(2)/(\pi/3)^2$ and $\cos(\pi/3)=\mu$. So \eqref{eq:projpcf} is not a fitted form: it is
``ratio, projected on $P$, in units of $\pi/3$'', and both constants come from
\eqref{eq:norms-derived}. \P{M7_oplus_formula}
\end{remark}

\begin{proposition}[Fibonacci sum as projection]\label{prop:oplus}
$F=P\oplus C:=\pi_{\PCF}(P,C,1)=\dfrac{(PC)\,\pi}{3\sqrt3}$.
\P{}
\end{proposition}

\begin{theorem}[Derivation of $\varepsilon_0$]\label{thm:eps0}
Apply \eqref{eq:projpcf} to the three data of the module direction---the modulus
$\mu=\sin(\pi/6)$ \eqref{eq:collapse}, the logarithmic generator $\ln\golden$ of the
multiplicative tower $\golden^{\sigma}$ of \eqref{eq:base}, and the half-turn $\pi$ accumulated
once around the fundamental domain (Proposition~\ref{prop:mobius-torus}):
\begin{equation}\label{eq:eps0}
\varepsilon_0=\pi_{\PCF}\!\bigl(\sin(\tfrac{\pi}{6}),\,\ln\golden,\,\pi\bigr)
=\frac{\ln\golden}{6\sqrt3}.
\end{equation}
\P{M8_epsilon0_from_projection}
\end{theorem}
\begin{proof}
By \eqref{eq:projpcf} with $a=\sin(\pi/6)=\tfrac12$, $b=\ln\golden$ and $c=\pi$,
\[
  \pi_{\PCF}\!\bigl(\tfrac12,\ln\golden,\pi\bigr)
  =\frac{\tfrac12\ln\golden}{\pi}\cdot\frac{1}{\sqrt3}\cdot\frac{\pi}{3}
  =\frac{\ln\golden}{6\sqrt3}.
\]
The three arguments are not free either: $\sin(\pi/6)$ is the modulus $\mu$ of
Theorem~\ref{thm:collapse} below, $\ln\golden$ is the logarithmic rate of the only tower the ring
\eqref{eq:base} generates, and $\pi$ is the monodromy of the fibre, of order two by
Proposition~\ref{prop:mobius-torus} below. The $\pi$ of the angular unit and the $\pi$ of the monodromy
cancel, which is why $\varepsilon_0$ is a pure multiple of $\ln\golden$ and why the scale it fixes
below satisfies $\varepsilon_0\Mpcf=\pi$ exactly.
\end{proof}

\begin{remark}[Torus, scale, and Sierpi\'nski attractor]
The microstate is realised on the torus
\begin{equation}\label{eq:torus}
T^2_{\PCF}=\C/\Lambda_{\PCF},\qquad \Lambda_{\PCF}=\Mpcf\,\Z[i],\qquad \tau=i ,
\end{equation}
whose scale is fixed by $\varepsilon_0$ of \eqref{eq:eps0}:
\begin{equation}\label{eq:Mpcf}
\Mpcf=\frac{6\sqrt3\,\pi}{\ln\golden}=\frac{\pi}{\varepsilon_0}\approx 67.846 .
\end{equation}
Equivalently, the \emph{Certainty Principle}
\begin{equation}\label{eq:certainty}
\varepsilon_0\,\Mpcf=\pi .
\end{equation}
The three eigenvalues $\lambda_k=\tfrac12\omega^{k}$ define an iterated function
system of three contractions of ratio $\tfrac12$,
\begin{equation}\label{eq:IFS}
T_k(z)=\tfrac12(z-v_k)+v_k,\qquad v_k=\tfrac12\omega^{k},\qquad k\in\{0,1,2\},
\end{equation}
whose Sierpi\'nski-type attractor has Hausdorff dimension
\begin{equation}\label{eq:hausdorff}
d_H=\frac{\log3}{\log2}\approx 1.585,\qquad 2^{d_H}=3 ,
\end{equation}
the Moran equation counting the three components $P,C,F$.
\end{remark}

% ---------------------------------------------------------------------

\subsection{The non-orientable fibre is independent of the torus}\label{ssec:mobius-torus}

The fibre, the isometry type of each map in the golden coupling, and the lattice scale that closes
the two.

The base torus of \eqref{eq:torus} carries a module direction and one non-orientable
identification of that direction, both fixed by data already in place. The module direction
is the phase $\Om(\tau)=\tfrac12\,e^{\,i\tau\ln\golden}$, of constant modulus
$|\Om(\tau)|=\tfrac12=\|\hat\Om\|$ (Proposition~\ref{prop:pcf-norms}).

\emph{Arithmetic type: an Eisenstein--Gauss transition.} The eigenvalue triangle
$\lambda_k=\tfrac12\omega^{k}$ ($\omega=e^{2\pi i/3}$) of Proposition~\ref{prop:pcf-norms}
sits in the Eisenstein lattice $\Z[\omega]$ (hexagonal, $120^{\circ}$); the torus modulus
$\tau=i$ sits in the Gauss lattice $\Z[i]$ (square, $90^{\circ}$). The golden coupling
$z=\golden y$ of \eqref{eq:base} carries the first into the second, and
Proposition~\ref{prop:coupling-isometries} identifies exactly which map is an isometry at each
step: the passage $120^{\circ}\to90^{\circ}$ is not a distortion of one lattice into another but a
change of the group that acts, and $\tau=i$ is the fixed point of the one that survives. The mediator is $\golden$ and no other: from \eqref{eq:base},
$\golden=1+1/\golden$, so its continued fraction is $[1;1,1,\dots]$ and $\golden$ is the
extremal irrational, most poorly approximated by rationals (Hurwitz~\cite{Hurwitz1891})---the
optimality that also places $\golden$ in phyllotaxis---and the modulus $\tfrac12$ is carried
uniformly on the fibre.

\begin{proposition}[The fibre is independent of the base algebra]\label{prop:mobius-torus}
The non-orientable identification lives in the fibre, not in the algebra of the torus, and is
therefore independent of whether that algebra commutes. Carried once around the fundamental
domain, the module direction accumulates the half-turn $\epsz\Mpcf=\pi$~\eqref{eq:certainty},
acting on the fibre of modulus $\tfrac12$ as $e^{i\pi}=-1$: the reflection $v\mapsto -v$, $\pi$
and not $2\pi$ because $\|\hat\Om\|=\tfrac12$, and orientation-reversing of order two,
$(-1)^2=1$. This datum is a $\Z/2$ reflection of the fibre; it carries no hypothesis on the
base commutator phase $UV=e^{2\pi i\theta}VU$. Hence the same fibre couples onto any torus---the
commutative torus ($\theta=0$) and the noncommutative torus $T^2_\theta$ that M-theory uses
($\theta$ irrational, \S\ref{sec:manin})---while preserving $|\Om|=\tfrac12$ on either.
\P{}
\end{proposition}
\begin{proof}
The fibre monodromy is $e^{i\pi}=-1$, of order two, $(-1)^2=1$, a $\Z/2$ datum that holds with
no condition on $\theta$ (\P{fibre\_independent\_of\_base}); the base commutes for
$\theta\in\Z$ and does not for $\theta$ irrational, the fibre being the same in both. The
modulus $\tfrac12$ is Proposition~\ref{prop:pcf-norms}, constant in $\tau$ since
$|\Om(\tau)|=\tfrac12$, and is preserved by the reflection.
\end{proof}



\begin{proposition}[$\Mpcf$ is fixed by the ring on the torus, not chosen]\label{prop:mpcf-forced}
The lattice scale of \eqref{eq:torus} is not a free parameter. Two data already in place fix it.
The module direction advances at the rate $\varepsilon_0$ of \eqref{eq:eps0} per unit of lattice
length, and carried once around the fundamental domain it must return the fibre to itself through
the order-two monodromy of Proposition~\ref{prop:mobius-torus}---the half-turn $\pi$, not $2\pi$,
because $\lVert\hat\Om\rVert=\tfrac12$. That closure condition,
\begin{equation}\label{eq:closure}
  \varepsilon_0\,\Mpcf=\pi ,
\end{equation}
has the unique solution $\Mpcf=\pi/\varepsilon_0$, and every factor of it is an invariant of
$\Rpcf$ on the torus:
\begin{equation}\label{eq:mpcf-invariants}
  \Mpcf\;=\;\frac{\pi}{\varepsilon_0}
  \;=\;\frac{3\pi}{\mu\,\lVert P\rVert\,\ln\golden}
  \;=\;\frac{d^{2}\,\pi}{\lVert v\rVert\,\ln\golden}
  \;=\;\frac{6\sqrt3\,\pi}{\ln\golden}\;\approx\;67.846 ,
\end{equation}
with $d=3$ the arity of \S\ref{ssec:arity}, $\mu=\tfrac12$ the modulus, $\lVert P\rVert=1/\sqrt3$
the projection norm and $\lVert v\rVert=\sqrt3/2$ the norm by which Lemma~\ref{lem:isometry}
divides the eigenvalue triad to make $V^{\dagger}V=\mathbf1$, and $\ln\golden$ the logarithmic rate
of the generator of \eqref{eq:base}. \P{M_PCF_eq_pi_div_eps0}
%
\P{M8\_epsilon0\_from\_projection, bulkBoundary\_isometry}%
\end{proposition}
\begin{proof}
The monodromy is $e^{i\theta}$ with $\theta=\varepsilon_0\Mpcf$ the angle accumulated over one
period. Proposition~\ref{prop:mobius-torus} requires that this act on the fibre of modulus
$\tfrac12$ as $v\mapsto-v$, of order two; on $[0,2\pi)$ the only such angle is $\theta=\pi$, whence
\eqref{eq:closure} and $\Mpcf=\pi/\varepsilon_0$, the solution being unique because
$\varepsilon_0>0$.

For \eqref{eq:mpcf-invariants}, \eqref{eq:projpcf} and \eqref{eq:eps0} give
$\varepsilon_0=\mu\lVert P\rVert\ln\golden/3$, so $\Mpcf=3\pi/(\mu\lVert P\rVert\ln\golden)$. The
isometric division of Lemma~\ref{lem:isometry} has $\lVert v\rVert=\sqrt3/2$, which is $\lVert
F\rVert$ of \eqref{eq:norms-derived}, and $\lVert P\rVert\lVert F\rVert=\tfrac12=\mu$ by
Proposition~\ref{prop:pcf-norms} with $\lVert C\rVert=1$. Hence
$\mu\lVert P\rVert=\lVert P\rVert^{2}\lVert v\rVert=\lVert v\rVert/d$, using
$\lVert P\rVert^{2}=1/3=1/d$, and $3\pi/(\lVert v\rVert\ln\golden/d)=d^{2}\pi/(\lVert
v\rVert\ln\golden)$. Substituting $\lVert v\rVert=\sqrt3/2$ and $d=3$ gives
$18\pi/(\sqrt3\ln\golden)=6\sqrt3\pi/\ln\golden$.
\end{proof}

\begin{remark}[What this rules out]\label{rmk:mpcf-not-free}
Had $\Mpcf$ been an independent scale, \eqref{eq:closure} would be a normalisation choice and the
Certainty Principle a convention. It is neither: the rate $\varepsilon_0$ is fixed by the ring's own
data (Theorem~\ref{thm:eps0}), the required monodromy is fixed by the fibre being of order two
(Proposition~\ref{prop:mobius-torus}), and the lattice scale is what closes the two. In particular
$\Mpcf$ is not tuned to the arity, the modulus or the isometric division---it is
\emph{their quotient}, \eqref{eq:mpcf-invariants}.
\end{remark}

\begin{proposition}[The golden coupling: which map is an isometry]\label{prop:coupling-isometries}
Four maps occur in the passage from the eigenvalue triangle to the torus, and the isometry type of
each is different and explicit.
\begin{enumerate}
\item[(i)] \emph{The Farish rotation is an ordinary isometry of $\R^{3}$.} The cyclic permutation
  $\rho(x,y,z)=(y,z,x)$ is orthogonal with $\det\rho=+1$ and $\rho^{3}=\mathrm{id}$, fixing the
  diagonal $a=\tfrac{1}{\sqrt3}(1,1,1)$. It preserves the inner product, hence all lengths and
  distances. The golden coupling is therefore not one plane but the $\rho$-orbit of \emph{three
  congruent planes}, congruent in the ordinary metric sense, with normals of equal length
  $s_\golden=\sqrt{1+\golden^{2}}=2\sin\tfrac{2\pi}{5}$ and equal mutual inclination
  $\arccos(-1/\sqrt5)$.
\item[(ii)] \emph{The Farish sense enters only in the projection.} Along $a$ the three axes are
  foreshortened equally, $\lVert\pi_a(e_k)\rVert=\sqrt{2/3}$, and appear at $120^{\circ}$. This is
  a property of how the figure is \emph{seen}, not of the objects, and it is the sense Farish fixed
  in 1822~\cite{fgc:farish22}.
\item[(iii)] \emph{On the torus the order three descends as an ordinary isometry too.}
  $\operatorname{Aut}(\Lambda_{\PCF})=\Z[i]^{\times}\cong\Z_4$ contains no element of order three,
  so the symmetry changes type rather than disappearing: it acts by \emph{translation} by a
  three-torsion point of $\mathbb{T}^{2}_{\PCF}[3]\cong(\Z/3)^{2}$. Translations are isometries of
  the flat torus, and since $3t=0$ gives $2t=-t$, the orbit $\{p,p+t,p+2t\}$ is
  \emph{equilateral in the ordinary metric}, of side $\lvert t\rvert$. Both $\Z_4$ (rotation,
  the Gauss face $\tau=i$) and $\Z_3$ (translation, the ternary face) therefore act by isometries.
\item[(iv)] \emph{Only the coupling map itself expands, and by an explicit factor.} The induced metric of
  $\iota(x+iy)=(x,y,\golden y)$ is
  \begin{equation}\label{eq:coupling-metric}
    \lVert\iota(x+iy)\rVert^{2}=x^{2}+(1+\golden^{2})y^{2},
  \end{equation}
  so $\iota$ preserves the real direction and expands the imaginary one by $s_\golden$. This does
  not contradict (i)--(iii): they concern different maps.
\end{enumerate}
Consequently the passage $120^{\circ}\leftrightarrow90^{\circ}$ is a change of the group under
which the structure is read---in $\R^{3}$ the projection along the Farish axis, on the torus
rotation versus translation---and not a morphism carrying one lattice onto another. \P{}
\P{farishRot\_cube, farishRot\_diag, farishRot\_preserves\_dot, farishRot\_preserves\_norm,
order\_three\_not\_in\_Z4, translation\_preserves\_difference, three\_torsion\_equilateral,
coupling\_metric, coupling\_expands\_imaginary, coupling\_preserves\_real}
\end{proposition}
\begin{proof}
(i) $\rho$ is a permutation matrix: $\rho^{\mathsf T}\rho=\mathbf 1$, $\det\rho=+1$,
$\rho^{3}=\mathrm{id}$, $\rho(1,1,1)=(1,1,1)$; and $\langle\rho u,\rho v\rangle=\langle u,v\rangle$
by inspection of the three terms. Congruence of the planes follows because $\rho$ is orthogonal, so
$\rho(n^{\perp})=(\rho n)^{\perp}$.
(ii) $\langle e_k,a\rangle=1/\sqrt3$ gives $\pi_a(e_k)=e_k-\tfrac13(1,1,1)$, of squared norm
$\tfrac23$, and $\langle\pi_a(e_j),\pi_a(e_k)\rangle=-\tfrac13$ for $j\neq k$, so
$\cos\angle=-\tfrac12$.
(iii) $3\nmid4$; $(p+t)-(q+t)=p-q$ in any abelian group, so translation preserves differences and
hence the quotient metric; and $3t=0\Rightarrow2t=-t$ makes the three pairwise differences $\pm t$.
(iv) is \eqref{eq:coupling-metric} evaluated at $x=1,y=0$ and at $x=0,y=1$.
\end{proof}


% ---------------------------------------------------------------------
\subsection{The single microstate, the meta-invariant, and the geometric
origin of $\tfrac12$ tied to the Euler product}\label{ssec:microstate}

The geometric origin of $\tfrac12$ via the angle $\pi/6$,
its anchoring to the critical line, and the factorial as connector.

\begin{theorem}[Trigonometric collapse of the norm product]\label{thm:collapse}
\begin{equation}\label{eq:collapse}
\|P\|\,\|C\|\,\|F\|
=\tan(\tfrac{\pi}{6})\cdot 1\cdot\cos(\tfrac{\pi}{6})
=\sin(\tfrac{\pi}{6})=\tfrac12 .
\end{equation}
\P{}
\end{theorem}

\begin{proposition}[Angle $\pi/3$ gives both invariants]\label{prop:angle}
$\sin(\tfrac{\pi}{6})=\cos(\tfrac{\pi}{3})=\mu$, and
\begin{equation}\label{eq:sigma-basel}
\sigma=\frac{\zeta(2)}{(\pi/3)^2}=\frac{\pi^2/6}{\pi^2/9}=\frac{9}{6}=\frac32 .
\end{equation}
\P{}
\end{proposition}

\begin{proposition}[Factorial as connector]\label{prop:connector}
$\dfrac{(\tfrac12)!}{\sqrt\pi}=\mu=\tfrac12$ \ (half face), and
$3!=6=|S_3|$ \ (integer face).
\P{}
\end{proposition}

\begin{remark}[$\tfrac12\leftrightarrow$ Euler product]\label{rmk:half-euler}
That the modulus $\tfrac12=\mu$ is anchored to the self-dual line is not an analogy: it is
Theorem~\ref{thm:funct-eq} together with Remark~\ref{rmk:half-selfdual}, where the fixed point of
$s\mapsto1-s$ is shown to be the same $\tfrac12$. The Euler product itself is
Proposition~\ref{prop:euler-product}, and it factors through the ternary level $\mu_3=\tfrac12$.
\end{remark}

% ---------------------------------------------------------------------
\subsection{Expansion of the spectrum mediated between $i$ and $\golden$}\label{ssec:spectrum}

Additive spectrum (pentagons) and multiplicative spectrum (Mersenne), joined by the $3/2$
structure; and the character $\chi_5$ read off the pentagon itself, with the Fibonacci criterion
that makes the classification of the primes effective.

\begin{equation}\label{eq:additive-mult}
\golden^n=F_n\golden+F_{n-1}\quad(\text{additive}),\qquad
3\,\golden^{\sigma(p)}=2^{p}\quad(\text{Mersenne}).
\end{equation}

\begin{proposition}[Mediated Mersenne]\label{prop:mersenne-med}
$3\,\golden^{\mu p\lng}=3\,(\sqrt2)^{p}$. \P{M12\_mersenne\_mediated}
\P{}
\end{proposition}

\begin{definition}[The quadratic character]\label{def:chi5}
$\chi_5$ is the Legendre symbol $\bigl(\tfrac{\cdot}{5}\bigr)$ on $\Z/5\Z$:
\begin{equation}\label{eq:chi5-values}
  \chi_5(0)=0,\qquad \chi_5(1)=\chi_5(4)=+1,\qquad \chi_5(2)=\chi_5(3)=-1,
\end{equation}
multiplicative on $(\Z/5\Z)^{\times}$, with $\sum_{a=0}^{4}\chi_5(a)=0$. \P{chi5, chi5_even, chi5_mul_on_units, chi5_sum_zero}

\end{definition}

\begin{proposition}[The pentagon dictates $\chi_5$]\label{prop:pentagon-chi5}
The four pentagonal cosines are exactly $\{\pm\golden,\pm\golden^{-1}\}$, so
\begin{equation}\label{eq:chi5-pentagon}
  \bigl|2\cos(\pi a/5)\bigr|=\golden^{\chi_5(a)},
  \qquad\text{equivalently}\qquad
  \chi_5(a)=\frac{\log\bigl|2\cos(\pi a/5)\bigr|}{\log\golden},
  \qquad a=1,2,3,4 .
\end{equation}
The character is therefore read off the pentagon itself, through
Theorem~\ref{thm:pentagon-id} ($\golden=2\cos(\pi/5)$), and needs no separate arithmetic input.
\P{norm_cosine_phi, phi_eq_two_cos_pi_fifth}
\end{proposition}
\begin{proof}
$2\cos(\pi/5)=\golden$ is Theorem~\ref{thm:pentagon-id}. The double-angle formula with
$\cos(\pi/5)=\golden/2$ and $\golden^{2}=\golden+1$ gives
$2\cos(2\pi/5)=\golden-1=\golden^{-1}$, and $\cos(\pi-x)=-\cos x$ carries the remaining two
cases: $\cos(3\pi/5)=-\cos(2\pi/5)$ and $\cos(4\pi/5)=-\cos(\pi/5)$. Taking absolute values gives
$\golden,\golden^{-1},\golden^{-1},\golden$ against $\chi_5=+1,-1,-1,+1$.
\end{proof}

\begin{proposition}[The logarithmic signature of $\golden$]\label{prop:log-signature}
\begin{equation}\label{eq:log-signature}
  \sum_{a=1}^{4}\chi_5(a)\,\log\bigl(2\sin(\pi a/5)\bigr)=-2\log\golden .
\end{equation}
The whole content is $\sin2x=2\sin x\cos x$: with $\sin(3\pi/5)=\sin(2\pi/5)$ and
$\sin(4\pi/5)=\sin(\pi/5)$ the sum collapses to $2\log\bigl(\sin(\pi/5)/\sin(2\pi/5)\bigr)$, and
$\sin(2\pi/5)/\sin(\pi/5)=2\cos(\pi/5)=\golden$. So the sine reads the same character the cosine
read in \eqref{eq:chi5-pentagon}, and $\golden$ is the value in both. \P{sin_ratio_eq_two_cos, pentagon_log_signature}

\end{proposition}

\begin{proposition}[Fibonacci criterion and the trichotomy]\label{prop:fib-criterion}
For every rational prime $q\neq2,5$,
\begin{equation}\label{eq:fib-criterion}
  F_q\;\equiv\;\Bigl(\frac{q}{5}\Bigr)\;=\;\chi_5(q)\pmod q ,
\end{equation}
and the value is determined by the residue class modulo $20$: $\chi_5(q)=+1$ (split) for
$q\equiv1,9,11,19$, $\chi_5(q)=-1$ (inert) for $q\equiv3,7,13,17$, and $\chi_5(5)=0$ (ramified),
$5$ being the only ramified prime by $\Delta_K=5$ \eqref{eq:trace-norm}. This is what makes
$\chi_5$ an \emph{effective} classification of the primes of $\Q(\sqrt5)$: the class of $q$ is
computed from the Fibonacci sequence of \eqref{eq:additive-mult}, without factoring anything in
$\OK$.  for the thirteen primes decided in Lean and the twenty-three verified numerically; \P{mersenne_bridge, phi_pow_fib}
the general congruence for every $q\neq2,5$ is Lucas' and enters as \Hyp{}
\end{proposition}

% ---------------------------------------------------------------------
\subsection{The rings as a tower (our reinterpretation of the $\Lambda$-structure of Borger)}\label{ssec:tower}

The multiplicative tower generated by $\golden$ carries a Frobenius lift for each prime, and the
lifts compose as the primes multiply. Both facts are established here.

\begin{definition}[The golden monoid and its Fibonacci coordinates]\label{def:golden-monoid}
Let $\langle\golden\rangle=\{\golden^{\,n}:n\in\Z\}$ be the multiplicative monoid generated by the
unit $\golden$ of \eqref{eq:base}, and let $F_n$ be the Fibonacci sequence, $F_0=0$, $F_1=1$,
$F_{n+2}=F_{n+1}+F_n$---the depth-two recurrence singled out in Theorem~\ref{thm:fib-min}. The
coordinates of the tower in the basis $\{\golden,1\}$ hold in the following generality: in
\emph{any} commutative ring containing an element $\alpha$ with $\alpha^{2}=\alpha+1$, induction
on $n$ gives
\begin{equation}\label{eq:binet}
  \alpha^{\,n+1}=F_{n+1}\,\alpha+F_n\qquad(n\ge0),
\end{equation}
of which the additive spectrum \eqref{eq:additive-mult} is the case $\alpha=\golden$.
Nothing of $\R$ or of $\sqrt5$ enters: only the quadratic relation and the recurrence. This is
why the same coordinates govern the tower in $\Rpcf$, in $\OK$, and modulo $q$---which is what
Proposition~\ref{prop:fib-criterion} exploits. \P{binet_general, phi_pow_fib}
\end{definition}

\begin{definition}[Frobenius lift on the golden monoid]\label{def:frobenius}
For each prime $p$ define
\begin{equation}\label{eq:frobenius-tower}
  \psi_p:\langle\golden\rangle\to\langle\golden\rangle,\qquad
  \psi_p(\golden^{\,n})=\golden^{\,pn},
  \qquad\text{so}\qquad
  \psi_p(\golden)=\golden^{\,p}=F_p\,\golden+F_{p-1}
\end{equation}
by \eqref{eq:binet}. \P{psiGolden, psi_on_powers, psi_golden_fib}
\end{definition}

\begin{remark}[Two maps, one symbol]\label{rmk:psi-two}
The $\psi_p$ of \eqref{eq:frobenius-tower} is \emph{not} the map $a\mapsto a^{p}\bmod p$ recalled in
\S\ref{sec:manin}, which is a ring endomorphism in the sense of Manin and Buium~\cite{manin13,buium05}.
Ours is exact rather than reduced modulo $p$, and it is an endomorphism of the multiplicative
\emph{monoid} $\langle\golden\rangle$, not of the ring $\Rpcf$: additivity is not preserved, since
$\psi_p(\golden+1)=\psi_p(\golden^{2})=\golden^{2p}\neq\golden^{p}+1$ for $p\ge2$. The restriction to
the multiplicative level is the point of the descent, not a limitation of it.
\end{remark}

\begin{proposition}[The lifts compose as the primes multiply]\label{prop:psi-functorial}
For all $p,q\ge1$ and all $x\in\langle\golden\rangle$,
\begin{equation}\label{eq:psi-functorial}
  \psi_p\circ\psi_q=\psi_{pq},\qquad \psi_1=\mathrm{id},
\end{equation}
so $p\mapsto\psi_p$ is a monoid homomorphism from $(\N_{\ge1},\times)$ into
$\mathrm{End}_{\mathrm{Mon}}\langle\golden\rangle$. This is the $\Lambda$-structure of
Borger~\cite{Borger} in its multiplicative form, and it is what makes the tower carry the
multiplicative structure of the Euler product \eqref{eq:euler-product}.
\P{psi_functorial, psi_one}
\end{proposition}
\begin{proof}
On a generator, $\psi_p(\psi_q(\golden^{\,n}))=\psi_p(\golden^{\,qn})=\golden^{\,pqn}=\psi_{pq}(\golden^{\,n})$,
and every element of $\langle\golden\rangle$ is such a generator; $\psi_1(\golden^{\,n})=\golden^{\,n}$.
\end{proof}

The remaining ingredient---that the Dedekind zeta of $\Q(\sqrt5)$ is the \emph{Euler colimit} of
the partial Euler products over finite sets of primes,
$\zeta_{\Q(\sqrt5)}(s)=\varinjlim_{S}E(S)(s)$---is Proposition~\ref{prop:euler-colimit} below, and it is what
Theorem~\ref{thm:places} uses on the finite side.

% ---------------------------------------------------------------------
\subsection{The zeta values, the Euler product, and the completed zeta}\label{ssec:zeta}

The zeta values, even and odd; the Dedekind zeta of $K$ with its local factors and its Euler
colimit; the Archimedean place; and the completed zeta as the product over all places, with its
functional equation. This is the first of two movements---building $\xi$; \S\ref{ssec:L1}
evaluates it at $s=1$. The bridge \eqref{eq:pi-arccos} ties $\pi$ to $\golden$ throughout.

\begin{equation}\label{eq:dedekind}
\zeta(s)=\frac{\zeta_{\Q(\sqrt5)}(s)}{L(s,\chi_5)},\qquad
L(1,\chi_5)=\frac{2\log\golden}{\sqrt5}.
\end{equation}
\P{eq:local-dedekind}%

\begin{theorem}[Basel value]\label{thm:basel}
$\zeta(2)=\pi^2/6$ (Euler). \P{M14\_basel}
\P{}
\end{theorem}

\begin{theorem}[General even zeta value (Euler)]\label{thm:even-zeta}
For every integer $k\ge1$,
\begin{equation}\label{eq:even-zeta}
\zeta(2k)=(-1)^{k+1}\frac{B_{2k}\,(2\pi)^{2k}}{2\,(2k)!},
\end{equation}
with $B_{2k}$ the Bernoulli numbers; $k=1$ recovers $\zeta(2)=\pi^2/6$
(Theorem~\ref{thm:basel}). \P{even_zeta_bernoulli} \N{thm:basel}
\end{theorem}
\begin{proof}
Euler's cotangent expansion
$\pi z\cot(\pi z)=1-2\sum_{k\ge1}\zeta(2k)\,z^{2k}$ \,($0<|z|<1$)\, and the
Bernoulli generating function
$\frac{w}{e^{w}-1}=\sum_{n\ge0}\frac{B_n}{n!}w^{n}$ are linked by
$\pi z\cot(\pi z)=\pi i z+\dfrac{2\pi i z}{e^{2\pi i z}-1}$. Substituting the
second series with $w=2\pi i z$ and using $(2i)^{2k}=(-1)^{k}4^{k}$, the
coefficient of $z^{2k}$ on the right is $(-1)^{k}B_{2k}(2\pi)^{2k}/(2k)!$;
matching it to $-2\zeta(2k)$ gives \eqref{eq:even-zeta}. The pentagonal,
$\chi_5$-twisted refinement, with generalised Bernoulli numbers
$B_{2k,\chi_5}$, extends this computation; it is taken up in
\S\ref{app:arithmetic} below and carried out independently in the $\F_1$-descent
companion~\cite{F1}.
\end{proof}

The remaining ingredients of $\xi$ are the product over primes and the factor that completes it.

\begin{proposition}[Euler product]\label{prop:euler-product}
For $\Re s>1$,
\begin{equation}\label{eq:euler-product}
  \sum_{n\ge1}n^{-s}=\zeta(s)=\prod_{p}\big(1-p^{-s}\big)^{-1},
\end{equation}
by unique factorisation in $\Z$. \P{regge_dirichlet_eq_zeta, regge_tower_is_euler_product}

\end{proposition}

\begin{definition}[The Dedekind zeta of $K$]\label{def:zetaK}
\begin{equation}\label{eq:zetaK}
  \zeta_K(s)\;=\;\sum_{\mathfrak a}N(\mathfrak a)^{-s}
  \;=\;\prod_{\mathfrak p}\bigl(1-N(\mathfrak p)^{-s}\bigr)^{-1},
  \qquad \Re s>1,
\end{equation}
the sum over integral ideals of $\OK=\Z[\golden]$ \eqref{eq:OK-vs-Rpcf} and the product over prime
ideals, $N$ the absolute norm. This is the object \eqref{eq:dedekind} factorises; it is defined
here rather than assumed. \P{idealNormK, idealEulerK, euler_from_ideals_K}

\end{definition}

\begin{proposition}[Splitting data and local factors]\label{prop:local-factors}
Every rational prime $p$ falls in one of the three classes of
Proposition~\ref{prop:fib-criterion}, with $\sum_i e_if_i=[K:\Q]=2$ in each:
\begin{equation}\label{eq:splitting}
  \begin{array}{lllll}
    \chi_5(p)=+1 & \text{split}    & g=2 & e=1,\ f=1 & N(\mathfrak p)=p\\
    \chi_5(p)=-1 & \text{inert}    & g=1 & e=1,\ f=2 & N(\mathfrak p)=p^{2}\\
    \chi_5(p)=0  & \text{ramified} & g=1 & e=2,\ f=1 & N(\mathfrak p)=p
  \end{array}
\end{equation}
Writing $u=p^{-s}$, the product over the prime ideals above $p$ is
\begin{equation}\label{eq:local-K}
  f_p(s)=\prod_{\mathfrak p\mid p}\bigl(1-N(\mathfrak p)^{-s}\bigr)^{-1}
  =\begin{cases}
     (1-u)^{-1} & \chi_5(p)=0,\\
     (1-u)^{-2} & \chi_5(p)=+1,\\
     (1-u^{2})^{-1} & \chi_5(p)=-1,
   \end{cases}
\end{equation}
The von Mangoldt function of $K$ weights a prime power by the log of its norm,
\begin{equation}\label{eq:LambdaK}
  \Lambda_K(\mathfrak p^{k})=\log N(\mathfrak p),
\end{equation}
so the weight is $\log p$ at a split prime and $2\log p$ at an inert one. In all three cases the
local factor factors as
\begin{equation}\label{eq:local-dedekind}
  f_p(s)=(1-u)^{-1}\cdot\bigl(1-\chi_5(p)u\bigr)^{-1},
\end{equation}
which is \eqref{eq:dedekind} place by place. The inert case is the identity $(1-u)(1+u)=1-u^{2}$.
\P{degree_formula_K, ideal_norm_values_K}
%
\P{euler\_from\_ideals\_K}%
\P{vonMangoldtK\_values}%
\P{local\_dedekind\_ramified\_K, local\_dedekind\_split\_K, local\_dedekind\_inert\_K}%
\end{proposition}

\begin{proposition}[The Euler product as a colimit]\label{prop:euler-colimit}
Let $\mathcal P$ be the poset of finite sets of primes under inclusion and
$E(S)(s)=\prod_{p\in S}f_p(s)$. Then
\begin{equation}\label{eq:colimit}
  \zeta_K(s)=\varinjlim_{S\in\mathcal P}E(S)(s),\qquad \Re s>1 .
\end{equation}
\emph{Existence} is coefficientwise stabilisation: the $n$-th coefficient is already determined
once $S$ contains the prime factors of $n$, and enlarging $S$ never alters it. \emph{Uniqueness} is
extensionality: a Dirichlet series is determined by its coefficients. This is the \emph{Euler}
colimit, not the framework colimit of \S\ref{ssec:web}.  (absolute convergence for \P{}
$\Re s>1$ enters as ) \Hyp{}
\end{proposition}

\begin{proposition}[The Archimedean local factor]\label{prop:archimedean}
The completing factor is the Mellin transform of the self-dual Gaussian
\eqref{eq:selfdual-gaussian}, i.e.\ the local zeta integral at the real place,
\begin{equation}\label{eq:gammaR-mellin}
  \Gamma_{\R}(s)\;:=\;\pi^{-s/2}\Gamma\!\left(\tfrac s2\right)
  \;=\;2\int_{0}^{\infty} g(x)\,x^{s-1}\,dx ,
\end{equation}
so it is produced by the object of \S\ref{ssec:origins} and not imported: the $\pi$ is
that of Lemma~\ref{lem:gamma-half} and the $\Gamma$ that of the Euler integral
\eqref{eq:gamma-int}. Summed over the integer direction of the Gauss lattice
$\Lambda_{\PCF}=\Mpcf\,\Z[i]$ of \eqref{eq:torus}, the same weight gives
\begin{equation}\label{eq:theta-lattice}
  \Theta(t)=\sum_{n\in\Z}e^{-\pi n^{2}t},\qquad
  \Theta(1/t)=\sqrt t\,\Theta(t)\quad(t>0),
\end{equation}
which is the involution $\tau\to-1/\tau$ of the torus \eqref{eq:torus} written in the Mellin
variable $t$, with unique fixed point $t=1$, i.e.\ $\tau=i$; \eqref{eq:theta-lattice} is
Jacobi's theta transformation~\cite{Jacobi1829}, and reading $\Gamma_{\R}$ as the local zeta
integral at the real place is Tate's~\cite{Tate1967}. The Boltzmann weight $e^{-nt}$ of the primon gas fails
\eqref{eq:theta-lattice} and is therefore excluded by that involution.
(Poisson summation enters as \Hyp{}) \P{gammaR_from_tate}
%
\P{theta\_from\_jacobi, theta\_fixed\_point\_unique}%
\end{proposition}
\begin{proof}
Substituting $u=\pi x^{2}$ in \eqref{eq:gammaR-mellin} gives
$2\int_{0}^{\infty}e^{-\pi x^{2}}x^{s-1}dx=\pi^{-s/2}\int_{0}^{\infty}e^{-u}u^{s/2-1}du
=\pi^{-s/2}\Gamma(s/2)$, which is \eqref{eq:gamma-int} read with $A=\pi$ and exponent
$s/2$. For \eqref{eq:theta-lattice}, Poisson summation applied to $f(x)=g(x\sqrt t)$,
whose transform is $\widehat f(\xi)=t^{-1/2}g(\xi/\sqrt t)$ by
Proposition~\ref{prop:selfdual-gaussian}, gives $\Theta(1/t)=\sqrt t\,\Theta(t)$
directly. The fixed point of $t\mapsto1/t$ on $t>0$ is $t=1$, since
$(t-1)(t+1)=0$ with $t+1>0$. The exclusion is a computation:
$\sum_{n\ge1}e^{-nt}=(e^{t}-1)^{-1}$ does not satisfy \eqref{eq:theta-lattice}.
\end{proof}

\begin{remark}[The value at the self-dual point]\label{rmk:eta-i}
The classical evaluation $\Theta(1)=\pi^{1/4}/\Gamma(\tfrac34)$ (Gauss;
see~\cite{Jacobi1829}) together with the reflection formula $\Gamma(\tfrac14)\Gamma(\tfrac34)=\pi\sqrt2$ gives
\begin{equation}\label{eq:eta-i}
  \Theta(1)=\sqrt2\,\eta(i)=\golden^{\mu\lng}\,\eta(i),\qquad
  \eta(i)=\frac{\Gamma(\tfrac14)}{2\,\pi^{3/4}},\qquad
  \Gamma_{\R}(\tfrac12)=2\sqrt\pi\,\eta(i),
\end{equation}
with $\golden^{\mu\lng}=\sqrt2$ of Proposition~\ref{prop:mediation}.  (the two classical \P{M4_sqrt2, M4_sqrt3}
evaluations enter as ) \Hyp{}
\end{remark}

\begin{theorem}[The partition is the product over all places]\label{thm:places}
Write the framework's spectral partition as the product of the Archimedean factor
\eqref{eq:gammaR-mellin} with the finite places, the latter being the Dirichlet
series \eqref{eq:euler-product}:
\begin{equation}\label{eq:places-product}
  \Xipcf(s)\;=\;\underbrace{\Gamma_{\R}(s)}_{\text{Archimedean}}\cdot
  \underbrace{\sum_{n\ge1}n^{-s}}_{\text{finite places}}
  \;=\;\Gamma_{\R}(s)\,\zeta(s)\;=\;\xi(s)\qquad(\Re s>1).
\end{equation}
Equivalently $\Xipcf(s)=\int_{0}^{\infty}t^{s/2-1}\,\omega(t)\,dt$ with
$\omega(t)=\tfrac12(\Theta(t)-1)$, so the partition is the Mellin transform of the
framework's own lattice sum. Reading $\xi$ as the product over \emph{all} places, the Archimedean one included, is the
formulation of L\'opez Pe\~na and Lorscheid~\cite{LopezPenaLorscheid}. The completed zeta is
therefore not posited: the finite places are the Euler product \eqref{eq:euler-product} and the Archimedean place is the self-dual
Gaussian (Proposition~\ref{prop:archimedean}); \S\ref{ssec:p-tower} reads the first as the Regge
tower and \S\ref{ssec:object} the second as the worldline weight.
\P{framework_partition, framework_partition_eq_completed_zeta, places_assembly}
\end{theorem}
\begin{proof}
The finite-place leg is Proposition~\ref{prop:euler-product}; the Archimedean leg is
Proposition~\ref{prop:archimedean}. Their product is $\pi^{-s/2}\Gamma(s/2)\zeta(s)$,
which is $\xi(s)$ by definition of the completed zeta. For the integral form, apply
\eqref{eq:gammaR-mellin} at each level $n$ with $A=\pi n^{2}$ and sum over $n\ge1$,
exchanging sum and integral by absolute convergence for $\Re s>1$.
\end{proof}

\begin{remark}[Consistency with the $\F_1$ companion]\label{rmk:places-F1}
$\Q(\sqrt5)$ is real quadratic, hence has two real places, and $\chi_5$ is even
($\chi_5(-1)=\chi_5(4)=+1$), so each real place contributes $\Gamma_{\R}$ and not the
odd factor. The Dedekind factorisation $\zeta_K=\zeta\cdot L(\cdot,\chi_5)$ of
Proposition~\ref{prop:local-factors} therefore distributes one $\Gamma_{\R}$ to each factor, and the one
attached to $\zeta$ is exactly \eqref{eq:gammaR-mellin}. The companion's Euler colimit
is indexed by finite sets of \emph{primes}; the Archimedean place is the one it leaves
out, absorbed there into the counting correction $\kappa_K$ of the Connes--Consani
type~\cite{ConnesConsaniSite}. Proposition~\ref{prop:archimedean} supplies it. \P{chi5_even}
\end{remark}


\begin{theorem}[Functional equation and the self-dual line]\label{thm:funct-eq}
The completed zeta,
\[
  \xi(s)=\pi^{-s/2}\Gamma(s/2)\zeta(s),
\]
satisfies
\begin{equation}\label{eq:funct-eq}
  \xi(1-s)=\xi(s),
\end{equation}
and the reflection $s\mapsto 1-s$ fixes the single point $s=\tfrac12$ of $\C$, while the line
\begin{equation}\label{eq:selfdual-line}
  \{s\in\C:\ \Re s=\Re(1-s)\}=\{s\in\C:\ \Re s=\tfrac12\}
\end{equation}
is the locus on which it preserves the real part. The point and the line are distinct statements
about the same involution, and it is the line that the cocone of \S\ref{ssec:commute} takes as a
leg.
functional\_equation\_fixed\_line, fixed\_line\_is\_facePhi, places\_assembly}
\end{theorem}
\begin{proof}
By Theorem~\ref{thm:places}, $\xi(s)=\int_{0}^{\infty}t^{s/2-1}\omega(t)\,dt$ for $\Re s>1$, with
$\omega(t)=\tfrac12(\Theta(t)-1)$. Splitting at $t=1$ and applying \eqref{eq:theta-lattice} on
$(0,1)$---Riemann's argument~\cite{Riemann1859}---,
\begin{equation}\label{eq:xi-split}
  \xi(s)=\frac{1}{s(s-1)}
        +\int_{1}^{\infty}\!\big(t^{s/2-1}+t^{(1-s)/2-1}\big)\,\omega(t)\,dt ,
\end{equation}
whose right-hand side is manifestly invariant under $s\leftrightarrow1-s$: the involution exchanges
the two exponents and fixes $1/(s(s-1))$. This gives \eqref{eq:funct-eq}, and by analytic
continuation for all $s$. The completing factor $\pi^{-s/2}\Gamma(s/2)$ is not imported: it is the
Archimedean local factor \eqref{eq:gammaR-mellin}, the $\pi$ of Lemma~\ref{lem:gamma-half} with the
$\Gamma$ of \eqref{eq:gamma-int}, and it carries at every $s$ the same $\pi$ that
Theorem~\ref{thm:basel} gives at $s=2$. The map $s\mapsto1-s$ is an involution of $\C$; its fixed
set is $\{s:2s=1\}$, i.e.\ the single value $s=\tfrac12$. For the line, $\Re(1-s)=1-\Re s$, so
$\Re s=\Re(1-s)$ holds precisely when $2\Re s=1$: this is \eqref{eq:selfdual-line}, and by
\eqref{eq:bridge} its abscissa is $\golden^{-\lng}$. The line is fixed \emph{pointwise} not by
$s\mapsto1-s$ but by the antiholomorphic $s\mapsto1-\bar s$; that reading is
\S\ref{ssec:extreme} below.
\end{proof}

\begin{remark}[The self-dual line is the modulus]\label{rmk:half-selfdual}
The $\tfrac12$ of \eqref{eq:funct-eq} is the same $\tfrac12$ that \S\ref{ssec:origins} fixed as
the unique fixed point of $x\mapsto1-x$ (Proposition~\ref{prop:phi-branch}) and as the midpoint
of the Galois pair (Proposition~\ref{prop:galois-seed}): the self-dual line of the completed
zeta \emph{is} the modulus $\mu=\tfrac12$, and by \eqref{eq:selfdual-line} its abscissa is
$\golden^{-\lng}$. Three involutions---$x\mapsto1-x$ on $\R$, the Galois conjugation of $K$, and
$s\mapsto1-s$ on $\C$---have here one fixed value; by Lemma~\ref{lem:galois-inv} the first two
are the same map.
\end{remark}

% ---------------------------------------------------------------------
\subsection{The fundamental formula and the entropy bridge}\label{ssec:L1}

With $\xi$ built, the second movement evaluates it at $s=1$: the class number, the residue, the
value $L(1,\chi_5)$, and the bridge that turns it into one bit.

\begin{proposition}[Class number one]\label{prop:class-number}
The Minkowski bound of $K$ is
\begin{equation}\label{eq:minkowski}
  M_K=\frac{n!}{n^{n}}\Bigl(\frac{4}{\pi}\Bigr)^{r_2}\sqrt{|\Delta_K|}=\frac{\sqrt5}{2}\approx1.11803,
  \qquad (n,r_1,r_2,\Delta_K)=(2,2,0,5),
\end{equation}
and $M_K<2$. Granting Minkowski's theorem---every ideal class contains an integral ideal of norm
at most $M_K$---that ideal has norm a positive integer below $2$, hence norm $1$, hence is the
unit ideal. Therefore
\begin{equation}\label{eq:hK}
  h_K=1 .
\end{equation}
(Minkowski's theorem enters as \Hyp{}) \P{minkowski_bound_lt_two}
%
\P{class\_number\_one\_K}%
\end{proposition}

\begin{theorem}[The fundamental formula]\label{thm:L1}
\begin{equation}\label{eq:L1}
  L(1,\chi_5)=\frac{2\log\golden}{\sqrt5}=\frac{T}{\pi\sqrt5}\approx0.4304089410 .
\end{equation}
Granting the Dirichlet class-number formula~\cite{dirichlet},
\begin{equation}\label{eq:cnf}
  \lim_{s\to1}(s-1)\zeta_K(s)=\frac{2^{r_1}(2\pi)^{r_2}h_KR_K}{w\sqrt{\Delta_K}}
  =\frac{4\cdot1\cdot\log\golden}{2\sqrt5}=\frac{2\log\golden}{\sqrt5},
\end{equation}
with $h_K=1$ \eqref{eq:hK}, $R_K=\log\golden$ \eqref{eq:regulator}, $w=2$, $\Delta_K=5$
\eqref{eq:trace-norm}, $r_1=2$, $r_2=0$; then \eqref{eq:dedekind} and
$\lim_{s\to1}(s-1)\zeta(s)=1$ give \eqref{eq:L1}. The arithmetic downstream of the formula is
proved; the formula itself is the input.  (the class-number formula enters as \Hyp{}) \P{L1_from_class_number_formula, L1_eq_period_over_pi_sqrt5}
%
\end{theorem}

\begin{theorem}[The entropy--$L$-value bridge]\label{thm:entropy-bridge}
With $\lambda_{\log}=\log2/\log\golden$, so that $\golden^{\lambda_{\log}}=2$ and
$\golden^{-\lambda_{\log}}=\tfrac12$,
\begin{equation}\label{eq:entropy-bridge}
  \frac{S_{\mathrm{BH}}}{k_B}
  =\lambda_{\log}\,R_K
  =\lambda_{\log}\cdot\frac{\sqrt5}{2}L(1,\chi_5)
  =\log 2 .
\end{equation}
The regulator cancels exactly: $\lambda_{\log}\log\golden=\log2$, and
$(\sqrt5/2)L(1,\chi_5)=\log\golden=R_K$ by \eqref{eq:L1}. The one bit per Planck area that
\S\ref{ssec:meta} and \S\ref{ssec:accum} carry is therefore not a coincidence with $\mu_3=\tfrac12$:
it is the regulator of $\Q(\sqrt5)$, measured in the unit $\lambda_{\log}$. \P{lambda_log_mul_regulator, sqrt5_half_L1_eq_regulator, entropy_bridge}

\end{theorem}

\begin{remark}[Three roles of one invariant]\label{rmk:regulator-roles}
$R_K=\log\golden$ occurs three times over: as the class-number residue \eqref{eq:cnf}, as the
exponent of the torus period $T=2\pi\log\golden$ \eqref{eq:regulator}, and through
\eqref{eq:entropy-bridge} as the entropy scale. It is also the rate out of which
$\varepsilon_0=\ln\golden/(6\sqrt3)$ is built \eqref{eq:eps0}. One arithmetic invariant of $K$,
four appearances.
\end{remark}

\begin{proposition}[$\tfrac12$ is where the packing is maximal]\label{prop:entropy-max}
For every $p\in(0,1)$ the binary entropy satisfies
\begin{equation}\label{eq:entropy-max}
  H(p)=-p\log_2p-(1-p)\log_2(1-p)\;\le\;1,
\end{equation}
with equality exactly at $p=\tfrac12$. The modulus is therefore not merely a point where the
entropy happens to equal one bit: it is the point where it is \emph{largest}. \P{entropy_slack_nonneg, binary_entropy_le_one, binary_entropy_half}

\end{proposition}
\begin{proof}
$\log x\le x-1$ for $x>0$ applied to $x=1/(2p)$ and $x=1/(2(1-p))$ gives
$p\bigl(-\log 2p\bigr)\le(1-2p)/2$ and $(1-p)\bigl(-\log 2(1-p)\bigr)\le(2p-1)/2$; the two
right-hand sides cancel, so $p\log 2p+(1-p)\log 2(1-p)\ge0$, which is \eqref{eq:entropy-max} after
dividing by $\log2$. Equality in $\log x\le x-1$ holds only at $x=1$, i.e.\ only at $p=\tfrac12$.
\end{proof}



\begin{theorem}[Odd zeta values over the tower]\label{thm:zeta-odd}
For $k\ge1$, $L(2k+1,\chi_5)\neq0$ and
\begin{equation}\label{eq:zeta-odd}
  \zeta(2k+1)=\frac{\zeta_K(2k+1)}{L(2k+1,\chi_5)} .
\end{equation}
Together with Theorem~\ref{thm:even-zeta} this closes the two families: $\pi$ produces the even
values, and $\golden$ through $\chi_5$ produces the odd ones, both over the same pentagonal points
$\{a/5\}$. It is what the loop hierarchy of \S\ref{ssec:p-tower} rests on, and it is derived here
rather than deferred.  (the global Dedekind factorisation enters as \Hyp{}) \P{odd_zeta_ratio}

\end{theorem}

% ---------------------------------------------------------------------
\subsection{The commutative diagram: convergence on $\mu$ and
$\sigma$}\label{ssec:commute}


The constant $\tfrac12$ is the apex of a cocone: eight
routes coincide on it. The argument is convergence, not a chain.

\[
\begin{tikzcd}[row sep=tiny, column sep=large]
\dfrac{(\tfrac12)!}{\sqrt\pi} \arrow[dddddr] & \\
\dfrac{\Gamma(\tfrac32)}{\Gamma(\tfrac12)} \arrow[ddddr] & \\
\|P\|\|C\|\|F\| \arrow[dddr] & \\
\cos(\tfrac{\pi}{3}) \arrow[ddr] & \\
\golden^{-\lng} \arrow[dr] & \\
\dfrac{\golden+\bar\golden}{2} \arrow[r] & \ \mu=\tfrac12 \\
\text{fix}(x{=}1{-}x) \arrow[ur] & \\
\text{fix}(\Re s{=}\Re(1{-}s)) \arrow[uur] &
\end{tikzcd}
\]

\begin{remark}[One principle, five fixed points]\label{rmk:pi-selfdual}
None of the last three legs is a separate mechanism. Each of the framework's constants is the fixed point of
an involution: $a=\pi$ of the Fourier transform (Proposition~\ref{prop:selfdual-gaussian}),
$\tau=i$ of $\tau\mapsto-1/\tau$ on the torus~\eqref{eq:torus}, $\tfrac12$ of the
Galois conjugation of $\Q(\sqrt5)$---whose fixed field is $\Q$ and whose midpoint is
\eqref{eq:galois-seed}---and $\tfrac12$ of both $x\mapsto1-x$
(Proposition~\ref{prop:phi-branch}) and
$s\mapsto1-s$ (Theorem~\ref{thm:funct-eq}). The cocone above collects the \emph{five} that land on the
modulus; the $\pi$ of \eqref{eq:selfdual-gaussian} is the same principle applied to the weight
rather than to the modulus. By Lemma~\ref{lem:galois-inv} the Galois conjugation and
$x\mapsto1-x$ are one map, so the five involutions are four distinct maps landing on one value.

A further symmetry is \emph{not} of this kind, and the contrast is informative: the order three of
the eigenvalue triad is not the fixed point of an involution but a symmetry that \emph{changes
type} on descending---a rotation of $\R^{3}$ above, a translation by a three-torsion point of the
torus below (Proposition~\ref{prop:coupling-isometries}). It is the same three that
Proposition~\ref{prop:triad-invariants} reads as the exponent of
$\prod_k|\lambda_k|=2^{-3}$. Five fixed points and one change of type.
\end{remark}

\begin{theorem}[Modulus diagram commutes]\label{thm:mu-diagram}
Six values and two identifications coincide on the apex:
\[
\dfrac{(\tfrac12)!}{\sqrt\pi}=\dfrac{\Gamma(\tfrac32)}{\Gamma(\tfrac12)}
=\|P\|\|C\|\|F\|=\cos(\tfrac{\pi}{3})=\golden^{-\lng}
=\dfrac{\golden+\bar\golden}{2}=\mu=\tfrac12 ,
\]
and $x=1-x\iff x=\mu$ on $\R$, $\Re s=\Re(1-s)\iff\Re s=\mu$ on $\C$.
The non-elementary legs are established at~\eqref{eq:half-factorial} (the Gamma value),
\eqref{eq:collapse} (the norm product), \eqref{eq:bridge} (the binary bridge
$\golden^{\lng}=2$, whence $\golden^{-\lng}=\tfrac12$) and~\eqref{eq:galois-seed} (the Galois
midpoint, by Lemma~\ref{lem:galois-inv}); the angle $\cos(\tfrac{\pi}{3})$
and the involution fixed point are elementary; the self-dual line of $\xi$ is
\eqref{eq:selfdual-line}, derived above rather than cited.

Two of the eight come from the same quadratic read on its two faces: the binary
$\golden^{-\lng}$ through \eqref{eq:bridge}, and the arithmetic
$\tfrac{\golden+\bar\golden}{2}$ through \eqref{eq:trace-norm}. That they agree is not a further
coincidence but the statement that the two readings of $\golden^{2}=\golden+1$ close on the same
value.
\P{}
\end{theorem}

\begin{remark}[The microstate is unique: no landscape to select]\label{rmk:no-landscape}
The eight routes of Theorem~\ref{thm:mu-diagram} leave no freedom in the modulus: the value $\mu=\tfrac12$ is
forced at once by the Gamma value, the norm product, the angle $\cos(\tfrac{\pi}{3})$, the binary bridge, the
involution fixed point and the self-dual line of the completed zeta. The microstate is therefore the \emph{unique} self-consistent solution, not one option
among many. This is what dissolves the vacuum-selection problem of \S\ref{sssec:ensemble}: there is no
$\sim 10^{500}$ landscape to choose from and no anthropic selection to invoke---the single microstate is fixed by
the convergence itself, and the abstract's ``unique self-consistent solution'' is established here. The uniqueness proved here is that of the modulus $\mu=\tfrac12$ as the algebraic fixed point of the eight coincident routes (Theorem~\ref{thm:mu-diagram}); its reading as the absence of a physical vacuum landscape is the interpretive step developed in \S\ref{sec:implications}, where the observer builds---rather than selects---the geometry.
\end{remark}

\[
\begin{tikzcd}[row sep=small, column sep=large]
\dfrac{\zeta(2)}{(\pi/3)^2} \arrow[dr] & \\
\dfrac{|{\rm rot}\,S_3|^2}{|S_3|} \arrow[r] & \ \sigma=\tfrac32
\end{tikzcd}
\]

\begin{lemma}[Orders in $S_3$ and the geometric leg]\label{lem:s3-orders}
The equilateral triangle of eigenvalues $\lambda_k=\tfrac12\omega^{k}$
(Proposition~\ref{prop:pcf-norms}) has symmetry group the symmetric group
$S_3\cong D_3$. Its order is $|S_3|=3!=6$, and its rotation subgroup is the
cyclic group $C_3=\{1,r,r^2\}$ of order $|{\rm rot}\,S_3|=3$. Hence
\begin{equation}\label{eq:sigma-geom}
\frac{|{\rm rot}\,S_3|^2}{|S_3|}=\frac{3^2}{6}=\frac{9}{6}=\tfrac32 ,
\end{equation}
so the geometric leg equals $\sigma$ rather than defining it.
\P{M11_factorial_six, s3_orders, sigma_geom_from_S3}
\end{lemma}

\begin{theorem}[Sum diagram commutes]\label{thm:sigma-diagram}
Both routes coincide on $\sigma=\tfrac32$: the analytic leg gives
$\dfrac{\zeta(2)}{(\pi/3)^2}=\dfrac{\pi^2/6}{\pi^2/9}=\dfrac{9}{6}$~\eqref{eq:sigma-basel}, and the
geometric leg gives $\dfrac{|{\rm rot}\,S_3|^2}{|S_3|}=\dfrac{9}{6}$
(Lemma~\ref{lem:s3-orders}); both equal $\sigma=\tfrac32$.
\P{}
\end{theorem}

\begin{theorem}[Methods master]\label{thm:master}
Both diagrams commute and the spectral invariants follow:
$\sigma+\mu=2$, $\sigma/\mu=3$; and the arithmetic they carry closes on the same modulus, the
completed zeta $\xi=\Gamma_{\R}\zeta$ of Theorem~\ref{thm:places} being self-dual exactly on
$\Re s=\mu$ (Theorem~\ref{thm:funct-eq}). \P{section2\_master, places\_assembly}
\P{fixed_line_is_facePhi, functional_equation_fixed_line, places_assembly, s_duality_exact, s_duality_fixed_point_unique, theta_from_jacobi}
\end{theorem}

\begin{remark}[The sum and the real places]\label{rmk:sum-places}
Recorded as an exact coincidence, with nothing asserted: the invariant $\sigma+\mu=2$ of
Theorem~\ref{thm:master} equals $r_1=2$, the number of real places of $K=\Q(\sqrt5)$---which is
what fixes the factor $\Gamma_\R^{2}$ in the completed zeta and the doubling
$\kappa_K=2\kappa_\Q$ of \S\ref{app:arithmetic}. No causal claim is attached and none is used
below.
\end{remark}


\begin{remark}[One bit, three readings of the same $\tfrac12$]\label{rmk:bit-three-readings}
Equation~\eqref{eq:entropy-bridge} is not an analogy between an arithmetic invariant and a physical
one; it is one quantity read three ways, and the three are the three things the framework has been
saying about $\tfrac12$ all along.

It is \emph{self-referential}. The regulator measures $K$ by its own fundamental unit: $R_K=\log\golden$
and $\golden$ generates the units of $\OK$ \eqref{eq:trace-norm}, so the field is measured in the
scale it itself produces. On the geometric side the same $\tfrac12$ is the fixed point of the
involution $s\mapsto1-s$ (Theorem~\ref{thm:funct-eq}), and of the involution $x\mapsto1-x$ that
Theorem~\ref{thm:mu-diagram} assembles with seven further routes to the same value---a fixed point,
which is the non-paradoxical self-reference of
\S\ref{sec:manin} and Definition~\ref{def:dsr}, where the recurrence is distributed at depth two and
so never depends directly on itself.

It is \emph{maximal}. By Proposition~\ref{prop:entropy-max}, $\tfrac12$ is where the binary entropy
attains its largest value: one bit is not one value among many but the most a single site can carry.
That is what makes the holographic bound a bound, and it is why $f_{\mathrm{crit}}=\tfrac12$ of
\S\ref{ssec:accum} is a threshold rather than a coincidence.

And it is \emph{arithmetic}. Through $\lambda_{\log}R_K=\log2$ that maximal bit \emph{is} the
regulator of $\Q(\sqrt5)$, measured in the logarithmic index. A quantity that refers to itself, that
is the most information a site can pack, and that is an invariant of a number field, are here the
same quantity; the framework does not relate them, it finds them to be one.
\end{remark}
% ---------------------------------------------------------------------
\subsection{The conductor, the scale, and the unit of spacing}\label{ssec:spacings}

Section~\ref{ssec:spectrum} fixed the residue class modulo $20$ that decides $\chi_5$
\eqref{eq:fib-criterion}. What that $20$ \emph{is} has not been said, and it is needed before any
statement about the spacings of zeros: a spacing is a number only once a unit is fixed, and if the
unit were free the statement would be about the unit rather than about the zeros.

\begin{proposition}[The conductor is derived]\label{prop:conductor}
The modulus $20$ is not chosen. It is the least common multiple of the two periods already in the
framework---the $4$ of the turn, $i^{4}=1$, and the $5$ of the pentagon
(Theorem~\ref{thm:pentagon-id})---and these are coprime, so the lcm is the product:
\begin{equation}\label{eq:conductor}
  \operatorname{lcm}(4,5)=20,\qquad \gcd(4,5)=1 .
\end{equation}
\P{conductor_is_derived}
\end{proposition}

\begin{proposition}[The two factors are one four and one five]\label{prop:four-cocone}
The $4$ of the turn and the $4$ that the binary register carries are not two numbers that agree.
They are one, and the identification is a reduction rather than an analogy. In $\F_5$ the element
$2$ satisfies the equation of the turn,
\begin{equation}\label{eq:two-is-i}
  2^{2}=-1 \ \text{ in } \F_5, \qquad i^{2}=-1 \ \text{ in } \C ,
\end{equation}
and the two cycles agree term by term: $2^{k}\bmod 5$ runs $2,4,3,1$ where $i^{k}$ runs
$i,-1,-i,1$. The reduction that carries one to the other exists because $5$ splits in the Gaussian
integers, $(2+i)(2-i)=5$; under $\Z[i]\to\Z[i]/(2-i)\cong\F_5$ the image of $i$ is $2$. What makes
this possible is $5\equiv1\pmod4$, so that $\lvert(\Z/5)^{\times}\rvert=4$, and $5$ is the only
prime with $p-1=4$. The $4$ of \eqref{eq:conductor} is therefore $\operatorname{ord}(i)$, the
conductor of $\Q(i)=\Q(\zeta_4)$, and the period of the binary register is that same $4$ seen
through the reduction---not a property of the factor $4$ itself, where $2^{m}\equiv0$ for
$m\geq2$ and there is no period at all.
\P{two_sq_eq_neg_one_mod_five, five_splits_gaussian, conductor_four_is_order_of_turn, binary_period_lives_in_the_five_component, four_cocone}
\end{proposition}

\begin{proposition}[Each factor governs one character, and only one]\label{prop:attribution}
The factorisation $20=4\cdot5$ separates the two characters exactly. Writing $r_4,r_5$ for the two
reductions, the pair is injective---the Chinese remainder theorem in the form needed here---and
each factor determines its own character while leaving the other free:
\begin{equation}\label{eq:attribution}
  r_4(a)=r_4(b)\ \wedge\ r_5(a)=r_5(b)\implies a=b,
\end{equation}
with $r_4$ determining $\chi_4$ and $r_5$ determining $\chi_5$, and neither determining the other.
The two witnesses are what keep this from being empty: $1$ and $13$ agree mod $4$ and differ in
$\chi_5$; $1$ and $11$ agree mod $5$ and differ in $\chi_4$. Consequently the classes
$\{1,9,11,19\}$ and $\{3,7,13,17\}$ of \S\ref{ssec:spectrum} are not tabulated lists but the fibres
of $r_5$, and the $\chi_5$ of \eqref{eq:chi5-pentagon} agrees with the quadratic fibre:
$\chi_5(n)=+1$ exactly when $n\equiv\pm1\pmod5$.
\P{conductor_splits, fiber_mod4, fiber_mod5, chi5_is_the_mod5_fiber, mod4_does_not_determine_chi5, mod5_does_not_determine_chi4, conductor_attribution}
\end{proposition}

\begin{proposition}[The scale is injective in the conductor]\label{prop:scale}
For $q,T$ with $2\pi<qT$ put $\mathrm{sc}(q,T)=2\pi/\log(qT/2\pi)$. Distinct conductors give
distinct scales at the same height:
\begin{equation}\label{eq:scale-injective}
  q_1\neq q_2 \implies \mathrm{sc}(q_1,T)\neq \mathrm{sc}(q_2,T).
\end{equation}
Hence ``the unfolded spacings have mean one'' selects one conductor: it is a statement about the
zeros, not about a choice of units. Without it the bound of
Definition~\ref{def:repulsion} would say nothing.
\P{spacingScale, scale_injective_in_conductor}
\end{proposition}

\begin{proposition}[The envelope carries both faces]\label{prop:envelope}
With $\mathrm{env}(q,T)=\frac{T}{2\pi}\log\frac{qT}{2\pi e}$, whose derivative is the density that
unfolds the ordinates,
\begin{equation}\label{eq:envelope-splits}
  \log\frac{qT}{2\pi e}=\underbrace{\log q}_{\text{field}}+\log T
  -\underbrace{\log 2\pi e}_{\text{universal}},
\end{equation}
and only the field part distinguishes: $q_1\neq q_2\Rightarrow\log q_1\neq\log q_2$. With $q=5$ the
first term is the $\golden$-face of \S\ref{ssec:commute} and the third the $\pi$-face. The counting
law $N(T)=\mathrm{env}(T)+S(T)$ and the unit $u_n$ it derives are asymptotic and enter as hypothesis
arguments \Conj; what is proved is that the conductor entering them is forced.
\P{envelope, envelope_splits, field_part_distinguishes, CountingLaw, DerivedUnit, the_unit_is_derived}
\end{proposition}

\begin{proposition}[The radial reading]\label{prop:li-modulus}
For $\rho\neq0$,
\begin{equation}\label{eq:li-modulus}
  \Bigl\lvert 1-\tfrac1\rho\Bigr\rvert=1 \iff \Re\rho=\tfrac12 :
\end{equation}
the critical line is exactly the preimage of the unit circle under $\rho\mapsto1-1/\rho$.
\P{li_modulus_on_line}
\end{proposition}

% ---------------------------------------------------------------------
\subsection{The extreme case, and the individual the ensemble replaces}\label{ssec:extreme}

The abstract began with an ensemble: the pair correlation of the zeros of $\zeta$ follows GUE
statistics, and so does the boundary side of low-dimensional holography. An ensemble is what one
writes for want of an individual. This section has produced the individual---the modulus
$\mu=\tfrac12$ of Theorem~\ref{thm:mu-diagram}. What remains is to say, with statements of our
own, what the ensemble was standing in for; and the answer is not a statistic but a pair.

\begin{definition}[The companion of the functional equation]\label{def:sdual-mate}
For $\rho\in\C$ set
\begin{equation}\label{eq:sdual-mate}
  \sigma\tau(\rho)=1-\bar\rho ,
\end{equation}
the reflection $s\mapsto1-s$ of Theorem~\ref{thm:funct-eq} composed with complex conjugation.
The two commute---$1$ is real---so $\sigma\tau$ is an involution, and it is
\emph{antiholomorphic}.
\P{sdualMate, sdual_conj_commute, sdualMate_involutive}
\end{definition}

\begin{proposition}[The self-dual line, as a fixed set]\label{prop:line-fixed}
$\sigma\tau$ preserves the ordinate and reflects the abscissa,
$\Im\sigma\tau(\rho)=\Im\rho$ and $\Re\sigma\tau(\rho)=1-\Re\rho$; hence
\begin{equation}\label{eq:line-fixed}
  \sigma\tau(\rho)=\rho \iff \Re\rho=\golden^{-\lng}.
\end{equation}
This is the sense in which the line \eqref{eq:selfdual-line} is \emph{self-dual}: not merely the
locus where $s\mapsto1-s$ preserves the real part, but the set fixed pointwise by an
involution---the antiholomorphic one. The holomorphic reflection fixes a single point of $\C$;
the line is what its composition with conjugation fixes. Its abscissa is the apex of
Theorem~\ref{thm:mu-diagram}.
\P{sdualMate_im, sdualMate_re, sdualMate_fixed_iff}
\end{proposition}

\begin{proposition}[The extreme case has arity two]\label{prop:arity-two}
If $\Re\rho\neq\golden^{-\lng}$ then $\sigma\tau(\rho)\neq\rho$ and
$\Im\sigma\tau(\rho)-\Im\rho=0$: off the line there are \emph{two} distinct points sharing one
ordinate, at separation exactly zero. The configuration that GUE statistics make improbable is
therefore not a statistic here but a concrete pair. The phenomenon the ensemble
describes---\emph{level repulsion}, the suppression of small spacings---is on this side a
statement of arity two, decided pointwise on one pair.
\P{extreme_case_is_binary}
\end{proposition}

\begin{proposition}[Angular and radial bound the same event]\label{prop:two-readings}
For $\rho\neq0$, by \eqref{eq:li-modulus},
\begin{equation}\label{eq:two-readings}
  \sigma\tau(\rho)\neq\rho \iff \bigl\lvert1-\tfrac1\rho\bigr\rvert\neq1 .
\end{equation}
The angular reading (arity two: a pair with zero separation) and the radial one (arity zero: a
modulus off unity) are two descriptions of one event. No hypothesis is needed: the radial reading
is Proposition~\ref{prop:li-modulus}.
\P{binary_iff_radial}
\end{proposition}

\begin{definition}[Measured repulsion]\label{def:repulsion}
Say $Z$ \emph{repels} at $m>0$ if any two distinct elements differ in ordinate by at least $m$, the
ordinates unfolded by the unit of \S\ref{ssec:spacings}:
\begin{equation}\label{eq:repulsion}
  \forall\,\rho\neq\tau\in Z:\quad \lvert\Im\tau-\Im\rho\rvert\ \geq\ m .
\end{equation}
This is \emph{level repulsion}: the suppression of small spacings that GUE statistics describe and
Poisson statistics do not. It is measured, not proved : over the $237$ spacings of $\gamma_1=14.13$ to \N{}
$\gamma_{238}=453.99$, unfolded by the unit of \S\ref{ssec:spacings}, the minimum is $0.2911$ and
none falls below $0.25$, let alone $0.10$, where independence would put $22.6$ and the Gaussian
ensemble $0.25$. The unfolded mean is $0.9979$ with $\zeta$'s own conductor $q=1$ against $1.8828$
with the foreign $q=20$, and every other conductor breaks it as well ($q{=}2\colon1.2026$,
$q{=}5\colon1.4733$)---which is Proposition~\ref{prop:scale} showing in the data: the scale is
injective in $q$, so the ordinates identify which conductor belongs to which $L$-function.
Only the \emph{sign} of $m$ enters anything below.
\end{definition}

\begin{figure}[htb]\centering
\includegraphics[width=\linewidth]{fig7_spacing_histogram.pdf}
\caption{Level repulsion in $\zeta$ zeros (Def.~\ref{def:repulsion}): (a) histogram of $237$ unfolded spacings ($\gamma_1{=}14.13$ to $\gamma_{238}{=}453.99$) against the GUE Wigner surmise and the Poisson exponential; (b) zoom on the repulsion region, with $\delta_{\min}{=}0.2911$ marked. The shaded zone $u{<}\delta_{\min}$ is where GUE suppresses spacings that Poisson permits.}
\label{fig:spacing-histogram}
\end{figure}

\begin{proposition}[Repulsion excludes the extreme case]\label{prop:repulsion-excludes}
If $Z$ repels at some $m>0$ then no two distinct elements share an ordinate; and since off the line
the pair $\{\rho,\sigma\tau(\rho)\}$ has separation exactly zero
(Proposition~\ref{prop:arity-two}), every $\rho\in Z$ closed under $\sigma\tau$ satisfies
$\Re\rho=\golden^{-\lng}$. The configuration the ensemble makes improbable and the one repulsion
excludes are the same configuration.
\P{repulsion_excludes_zero_spacing, repulsion_excludes_the_extreme_case}
\end{proposition}

\begin{theorem}[The weak form, and what would suffice]\label{thm:weak-form}
Closure under $\sigma\tau$ alone gives
\begin{equation}\label{eq:E-iff-no-sharing}
  \bigl(\forall\rho\in Z:\ \Re\rho=\golden^{-\lng}\bigr)
  \iff \text{no two distinct elements of } Z \text{ share an ordinate}.
\end{equation}
The right-hand side is strictly weaker than Definition~\ref{def:repulsion}: repulsion implies it,
not conversely. It is what any future argument would discharge in place of the metric bound.
\P{E_iff_no_shared_ordinate}
\end{theorem}

\begin{proposition}[Repulsion gives the modulus]\label{prop:repulsion-modulus}
Under the same hypotheses, $\lvert1-1/\rho\rvert=1$ exactly. The measured repulsion does not merely
constrain the modulus: it \emph{produces} it.
\P{repulsion_gives_the_modulus}
\end{proposition}

\begin{theorem}[Zeros on the apex, with its inputs in view]\label{thm:zeros-apex}
Let $Z\subset\C$ be closed under $\sigma\tau$ and repel at some $m>0$
(Definition~\ref{def:repulsion}). Then every $\rho\in Z$ satisfies $\Re\rho=\golden^{-\lng}$, the
apex of Theorem~\ref{thm:mu-diagram}.

Two inputs are open and are named here rather than hidden. Closure needs the conjugation half of
\eqref{eq:sdual-mate}, $\xi(\bar s)=\overline{\xi(s)}$, not established in this paper \Conj; and
the repulsion is measured . Everything between them is proved, and of the measurement only \N{eq:sdual-mate}
the \emph{sign} enters.
\end{theorem}

\begin{remark}[What the argument does not use]\label{rmk:no-statistics}
Theorem~\ref{thm:zeros-apex} uses the geometry of $\C$ and the sign of one bound. It uses no
joint density, no sine kernel, no form factor, no asymptotic statistics of the spacings. That is
the point of stating the extreme case at arity two: the exclusion the criterion needs is about
\emph{one pair}, and a pair is an individual. Where the ensemble was standing in for an object,
the object here is the modulus, and what the statistics described is a consequence of its
geometry rather than a substitute for it. The chain that discharges the two open inputs, and the
measured bound itself, belong to the companion papers~\cite{HP,F1}; note that they are stated
there for $L(s,\chi_5)$, whose functional equation is a different theorem from
\eqref{eq:funct-eq}, proved separately in~\cite{F1}.
\end{remark}


\begin{definition}[The ensemble's own object]\label{def:pair-correlation}
The statistics invoked in \S\ref{sec:intro} are those of the GUE pair-correlation kernel
\begin{equation}\label{eq:sine-kernel}
  K(u)=1-\Bigl(\frac{\sin\pi u}{\pi u}\Bigr)^{2} .
\end{equation}
It is written here---the attribution is at \S\ref{sec:intro}---so that the object the paper opened
with is a node of this development rather than a gap in it. It vanishes as $u\to0$, which is the
suppression of small spacings that Definition~\ref{def:repulsion} measures, and tends to $1$ for
large $u$, which is decorrelation. Nothing below uses it. \Conj
\end{definition}

\begin{figure}[htb]\centering
\includegraphics[width=\linewidth]{fig8_sine_kernel.pdf}
\caption{The GUE pair-correlation kernel $K(u)=1-(\sin\pi u/\pi u)^{2}$ of Def.~\ref{def:pair-correlation}. The red point at $K(0){=}0$ is level repulsion---the suppression of small spacings that Def.~\ref{def:repulsion} measures. Blue points mark $K(n){=}0$ for $n\in\Z\setminus\{0\}$ (no long-range correlation). The kernel tends to $1$ for large $u$ (decorrelation).}
\label{fig:sine-kernel}
\end{figure}

\begin{remark}[Envelope yes, fluctuations no]\label{rmk:envelope-not-fluctuations}
The ensemble stands on a distinction between the smooth spectral density and the fluctuations
around it. A deterministic operator built without reference to $\zeta$ determines the
envelope---the global distribution of zeros---but the pair-correlation statistics reside in the
fluctuations, and reproducing them would require either an exact spectrum or a dynamical mechanism
generating the level repulsion of Definition~\ref{def:repulsion}. \emph{This is intrinsic to the
construction, not a gap in the work}: it is why the present subsection proves a statement about
\emph{one pair} rather than about a distribution. The joint density and the form factor are
declared in the development and used in no signature; that is the check that nothing here rests on
them. \Empty
\end{remark}

\begin{remark}[A geometric origin for the constraint]\label{rmk:geometric-origin}
What the framework does supply is a candidate \emph{origin} for the constraint that makes such
statistics possible, and it is already proved above: the ratio $\sigma/\mu$ is the arity $d=3$
(Remark~\ref{rmk:spectral-origin}), the modulus contracts strictly, $\|\hat\Omega\|=\tfrac12<1$
(Remark~\ref{rmk:no-diagonal}), and the same $3$ is the exponent of
$\prod_k\lvert\lambda_k\rvert=2^{-3}$ in Proposition~\ref{prop:triad-invariants}. The reading that
eigenvalue correlations in this regime arise from \emph{geometric constraint}---the no-diagonal
condition---rather than from dynamical instability is advanced in the companion paper~\cite{HP};
it is a proposal there and not a theorem here, and is recorded as such. What is established in
this section is the regime, not the reading. \Open
\end{remark}

\begin{remark}[Not the tower spacing]\label{rmk:not-the-tower-spacing}
One caution, now that spacings are in play: the tower has its own spacing, the regulator
$R_K=\log\golden$ of \eqref{eq:regulator}, numerically $0.481212$, while the unit that unfolds
ordinates at $T=100$ is $1.435589$ (Proposition~\ref{prop:scale}). They share the symbol $\golden$
and are different objects.
\P{tower_spacing_is_the_regulator}
\end{remark}

% ---------------------------------------------------------------------
\subsection{Formal verification}\label{ssec:status}

Each result of this section carries, in brackets, the name of its Lean~4
theorem. The development that accompanies the paper is the self-contained
file \texttt{CW6\_complete\_v2.lean}. It compiles under \texttt{leanprover/lean4}~v4.29.0 with
Mathlib~v4.29.0: \texttt{lake build} completes with no errors and no warnings, and the file
declares \emph{no} \texttt{axiom} and contains no \texttt{sorry}. The file closes with the
framework colimit of \S\ref{ssec:web} below (namespace \texttt{PCFColimit}): the frameworks the paper
treats---Polyakov, Maldacena, ER$=$EPR, Huerta--Schreiber, Kiely---as partial spectral objects
together with the PCF core, whose fields are populated by theorems of this same file
(\texttt{modulus\_Omega}, \texttt{holographic\_area}, \texttt{kk\_at\_PCF},
\texttt{two\_pow\_hausdorff}, the P/C/F norms), the cocone into the vertex $T_{\PCF}$
(\texttt{cocone\_property}), and the universal property that makes the vertex unique
(\texttt{colimit\_universal}, \texttt{framework\_colimit}).
The external results the file uses
enter as explicit hypothesis arguments in the types of the theorems that consume them---where
the reader can see them---rather than as axioms; the classical analytic inputs
(\texttt{gammaR\_from\_tate}, \texttt{theta\_from\_jacobi}) and the combinatorial ones
(\texttt{ladder\_from\_scott}, \texttt{e8\_dim\_from\_kissing}, \texttt{positroid\_from\_kp},
\texttt{no\_separable\_parts}) are carried that way. Two further inputs enter the same way in
\S\ref{ssec:extreme} and are recorded here because they are the only open ones in that
subsection: \texttt{XiConjClosed}, the closure $\xi(\bar s)=\overline{\xi(s)}$ of the zeros
under conjugation, which is \emph{not} established in this paper \Conj; and
\texttt{MinSpacing}, a positive minimal ordinate separation, which is measured rather than
proved . Of the second only its sign is used. Neither is an axiom: both are hypothesis \N{thm:zeros-apex}
arguments in the type of Theorem~\ref{thm:zeros-apex}. One decision procedure uses
\texttt{native\_decide} rather than \texttt{decide}---the sweep of $16\times17\times18$ triples behind
Proposition~\ref{prop:interval-uniqueness}---which extends the trusted base to Lean's compiler
as well as its type kernel; every other proof rests on the kernel alone. The golden ratio is
defined, $\golden:=(1+\sqrt5)/2$, and the two facts that anchor the
construction are theorems proved from that definition, not assumptions. The
relation $\golden^2=\golden+1$ (\texttt{phi\_sq}) is derived from the
definition and is~\eqref{eq:phisq}. The modulus $\|\hat\Omega\|=\tfrac12$
(\texttt{modulus\_Omega}) is likewise a theorem. Its value is reached by the routes assembled in
Theorem~\ref{thm:mu-diagram}, its component structure is
Proposition~\ref{prop:pcf-norms}, and the trigonometric collapse of the norm product is
Theorem~\ref{thm:collapse}; the three quantities that coincide there and cease to coincide
afterwards are separated in Proposition~\ref{prop:triad-invariants}.

\smallskip
\sloppy\noindent\textbf{Compilation note (single point of record).} The development declares no
axiom: every classical input is a hypothesis argument in the type of the theorem that consumes it.
\texttt{lake build} completes against Mathlib, confirming the spelling
of every library lemma the proofs call. The lemma names, grouped by what they serve,
are recorded here so that a future Mathlib rename breaks the build visibly rather than
passing silently:
\begin{itemize}\itemsep1pt\parskip0pt
\item \emph{Completed zeta and the places} (\S\ref{ssec:zeta}): \texttt{Gamma}$_\R$,
  \texttt{Gamma}$_\R$\texttt{\_ne\_zero\_of\_re\_pos}, \texttt{riemannZeta\_def\_of\_ne\_zero},
  \texttt{zeta\_eq\_tsum\_one\_div\_nat\_cpow}, \texttt{riemannZeta\_eulerProduct\_tprod},
  \texttt{completedRiemannZeta\_one\_sub}, \texttt{integral\_gaussian}, \texttt{Real.sq\_sqrt},
  \texttt{Real.Gamma\_eq\_integral}.
\item \emph{The generator and its conjugate} (\S\ref{ssec:generator}), \emph{the arithmetic
  seed} (\S\ref{ssec:origins}) \emph{and the cocone} (\S\ref{ssec:commute}):
  \texttt{Real.sq\_sqrt}, \texttt{Real.Gamma\_add\_one}, \texttt{Real.Gamma\_pos\_of\_pos}.
\item \emph{The eigenvalue triad} (Proposition~\ref{prop:triad-invariants}):
  \texttt{Complex.exp\_ofReal\_mul\_I\_re}, \texttt{Complex.exp\_nat\_mul},
  \texttt{Complex.exp\_two\_pi\_mul\_I}, \texttt{Real.cos\_pi\_sub},
  \texttt{Real.rpow\_mul}, \texttt{Real.rpow\_natCast}, \texttt{Real.rpow\_neg}.
\item \emph{The companion and the extreme case} (\S\ref{ssec:extreme}): \texttt{map\_sub},
  \texttt{map\_one} on \texttt{starRingEnd}, \texttt{Complex.conj\_re},
  \texttt{Complex.conj\_im}, \texttt{Complex.ext\_iff}.
\item \emph{Theta and Poisson} (\eqref{eq:theta-lattice}):
  the file defines \texttt{Theta} directly without Mathlib's \texttt{jacobiTheta}.
\item \emph{Pentagon and Fibonacci} (\S\ref{ssec:spectrum}, \S\ref{ssec:tower}):
  \texttt{Real.cos\_pi\_sub}, \texttt{Real.cos\_two\_mul}, \texttt{Nat.fib}.
  The general form of \eqref{eq:fib-criterion} would need Frobenius in \texttt{ZMod q}; it is
  stated here by cases.
\item \emph{Even zeta values} (Theorem~\ref{thm:even-zeta}): the name of the even-value lemma,
  and \textbf{whether it is stated with \texttt{bernoulli} or with \texttt{bernoulli'}}. The two
  differ only at $B_1$ and agree at every even index $\ge2$, so the statement is the same either
  way; the identifier is not.
\item \emph{Group orders} (Lemma~\ref{lem:s3-orders}): the index-two fact for $A_n$,
  \texttt{two\_mul\_card\_alternatingGroup} or its equivalent.
\item \emph{Matrices} (Theorem~\ref{thm:faces}): \texttt{Matrix.mul\_inv\_rev},
  \texttt{Matrix.transpose\_nonsing\_inv}.
\item \emph{Earlier blocks}: \texttt{Gamma\_one\_half\_eq}, \texttt{riemannZeta\_two},
  \texttt{Real.sqrt\_eq\_rpow}, \texttt{rpow\_natCast}, \texttt{Gamma\_add\_one},
  \texttt{arccos\_cos}, and for the M6 uniqueness \texttt{Nat.Prime.irrational\_sqrt},
  \texttt{Irrational.intCast\_mul}, \texttt{Int.not\_irrational}.
\end{itemize}
\noindent Since \texttt{lake build} completes, any renaming or removal of these lemmas in a
future Mathlib breaks the build rather than passing silently. The numerical companion is the
pytest suite in \PytestInline{tests/} (run via \texttt{pnpm run verify:py}),
whose tests are keyed by the \TeX{} labels of this file: of the $162$ labelled equations,
$148$ carry a Lean proof and $119$ a numerical check, $102$ carrying both---the tests evaluating
identities at up to $25$-digit precision and, where a claim is a correspondence with measurement,
comparing against the measured value. The eight carrying neither a proof nor a check of their
own are the five correspondence statements \eqref{eq:bridge-S}--\eqref{eq:bridge-ER} plus
\eqref{eq:dS-swampland}; \eqref{eq:cnf} whose Lean form is the hypothesis
\texttt{hCNF} of \texttt{L1\_from\_class\_number\_formula};
the Polyakov action \eqref{eq:polyakov} (an  input); and the \Hyp{}
Lawvere--Yanofsky diagonal \eqref{eq:LY}, which states the problem rather than a result.
The global Dedekind factorisation \eqref{eq:dedekind} is counted among the $147$: its Lean
backing sits at the local factors \eqref{eq:local-dedekind} that compose it. The figures are generated by
\texttt{CW6\_all\_figures\_v2.py} and their assertions checked, separately and against the labels of
this file, by \texttt{CW6\_figures\_verify\_v2.py}.


% ===================================================================
